\documentclass[11pt]{article}
\usepackage[a4paper,top=2.4cm,bottom=2.5cm,left=2.0cm,right=2.0cm]{geometry}
\usepackage{amsmath,amssymb}
\usepackage[UTF8]{ctex}
\usepackage{fontspec}
\setmainfont{TeX Gyre Termes}
\usepackage{booktabs}
\usepackage{graphicx}
\usepackage{pdflscape}
\usepackage[section]{placeins}
\usepackage{float}
\graphicspath{{./}}
\renewcommand{\topfraction}{0.92}
\renewcommand{\bottomfraction}{0.8}
\renewcommand{\textfraction}{0.07}
\renewcommand{\floatpagefraction}{0.75}
\usepackage{xcolor}
\definecolor{revisionblue}{RGB}{0,76,153}
\newcommand{\revmark}[1]{\textcolor{revisionblue}{\textbf{[R#1]}}\ }
\newcommand{\revdeleted}[1]{\par\noindent\textcolor{revisionblue}{\small\itshape [#1]}\par}
\newcommand{\revblue}[1]{\textcolor{revisionblue}{#1}}
\usepackage{tikz}
\usetikzlibrary{arrows.meta,positioning}
\usepackage{caption}
\captionsetup{font=small,labelfont=bf}
\usepackage{titlesec}
\titleformat{\section}{\normalfont\large\bfseries}{\thesection}{0.6em}{}
\titleformat{\subsection}{\normalfont\normalsize\bfseries}{\thesubsection}{0.5em}{}
\titleformat{\subsubsection}{\normalfont\normalsize\bfseries}{\thesubsubsection}{0.5em}{}
\titlespacing*{\section}{0pt}{1.3ex plus .2ex minus .2ex}{0.6ex}
\titlespacing*{\subsection}{0pt}{1.0ex plus .2ex minus .2ex}{0.5ex}
\usepackage{fancyhdr}
\pagestyle{fancy}
\fancyhf{}
\fancyhead[L]{\small\itshape Experimental Astronomy}
\fancyhead[R]{\small\thepage}
\renewcommand{\headrulewidth}{0.4pt}
\usepackage{cite}
\usepackage{hyperref}
\hypersetup{hidelinks}
\setlength{\emergencystretch}{3em}
\newcommand{\keV}{\,\mathrm{keV}}
\newcommand{\MeV}{\,\mathrm{MeV}}
\newcommand{\cms}{\,\mathrm{cm^{-2}\,s^{-1}}}
\newcommand{\phcms}{\,\mathrm{ph\,cm^{-2}\,s^{-1}}}
\newcommand{\cps}{\,\mathrm{s^{-1}}}
\newcommand{\wii}{W_{511}}
\newcommand{\aeff}{A_{\mathrm{eff}}}
\newcommand{\fref}{F_{\mathrm{ref}}}
\newcommand{\fthree}{F_{3\sigma}}

\begin{document}
\thispagestyle{fancy}
\begin{flushleft}
{\LARGE\bfseries 气球载 511 keV Laue 透镜 TES 望远镜本底和未分辨线灵敏度的探测器耦合 Monte Carlo 估计 \par}
\vspace{1.1em}
{\large 作者信息暂隐\par}
\vspace{0.3em}
{\normalsize\itshape 单位信息暂隐\par}
\end{flushleft}
\vspace{0.7em}
\noindent\rule{\textwidth}{0.4pt}
\vspace{0.5em}

\noindent\textbf{摘要}\hspace{0.7em}
利用气球载 Laue 透镜望远镜探测紧致 511 keV 辐射，必须在共同探测器响应与事例选择框架下同时评估聚焦信号、瞬发本底和活化本底。本文采用端到端 Geant4/MEGAlib 链，将聚焦天体信号、\revblue{八类瞬发大气粒子族}及按产生位置抽样的活化衰变分别输运至两套完整探测器--低温系统质量模型。每套设计内的全部事例流统一接受 $510.58$--$511.42\keV$ 科学窗、公共时间事例构建、探测器响应、主动反符合及拓扑/孔径选择。质量模型 A 第 15 天直接选后本底率为 $5.3057\times10^{-2}\cps$。对其带权选后延迟率作两个彼此独立的投影后，稀释制冷机混合室及各级冷端区域占 $32.18\%$，而铜材料占 $70.87\%$。这一来源诊断促使本文把质量模型 B 作为完整备选构型进行检验。质量模型 B 的选后有效面积为 $15.08544\pm0.04503\,\mathrm{cm^2}$，保留质量模型 A 的 $99.75\%$，同时把第 15 天直接选后本底率降至 $1.0747\times10^{-2}\cps$。对于由 IBIS 上限换算并定义在大气层外的唯一 $3\sigma$ 等效参考通量 $\fref=2.4\times10^{-4}\phcms$，20 d、81 节点折叠得到质量模型 B 的信号计数 4028.15、本底计数 18568.89，计数显著度为 29.561；相应的大气层顶 $3\sigma$ 阈值为 $(2.436\pm0.223)\times10^{-5}\phcms$。在给定源项、响应、计时及选择条件下，质量模型 B 将 20 d 本底降至质量模型 A 的 $21.19\%$，并使流量阈值改善 2.227 倍；若要把这一改善归因于某个单独设计变化，仍需专门的消融研究。

探索性的两流重加权还将安静 LEO、月壤辐射场和日--地 L2 识别为三类不同的后续环境问题；所得筛查指数不是环境匹配的任务灵敏度预报。

\vspace{0.7em}
\noindent\textbf{关键词}\hspace{0.7em}伽马射线仪器 $\cdot$ 511 keV 湮灭线 $\cdot$ 过渡边缘传感器微量热计 $\cdot$ Laue 透镜 $\cdot$ 气球载望远镜 $\cdot$ Monte Carlo 本底模拟

\vspace{0.3em}
\noindent\rule{\textwidth}{0.4pt}
\vspace{1.0em}

\section{引言}
\label{sec:intro}

银河系 511 keV 湮灭辐射是研究银河系正电子产生、传播和湮灭环境的重要观测探针。其谱学特征包括接近 511 keV 的双光子谱线和位于较低能量处的正电子素连续谱。正电子可与电子直接湮灭，或先形成对位正电子素并产生两个约 511 keV 光子；正交正电子素则通过三光子衰变形成连续谱。因此，谱线强度、轮廓和正电子素连续谱可以约束正电子减速和湮灭所在介质的性质 \cite{Prantzos2011,Jean2006}。在近似稳态假设下，INTEGRAL/SPI 测得的银河系 511 keV 总通量对应约几 $\times10^{43}\,\mathrm{e^+\,s^{-1}}$ 的正电子湮灭率 \cite{Prantzos2011,Kierans2020}。候选正电子来源包括恒星爆发中合成的 $\beta^+$ 不稳定核素、能够产生正负电子对的致密天体，以及更具推测性的粒子物理情景，例如由轻质量暗物质湮灭或衰变注入低能正电子。现有观测尚不能唯一确定其中的主导贡献 \cite{Prantzos2011,Siegert2016AA,Siegert2016V404}。

过去的宽视场观测已经建立了 511 keV 辐射的弥散形态。INTEGRAL/SPI 显示出明亮的中心核球成分和较弱的银盘成分，20 年数据分析进一步重建出中心、宽核球与银盘结构，但额外关联结构的证据仍较边际 \cite{Knodlseder2005AllSky,Siegert2016AA,Yoneda2025}。COSI 2016 年气球飞行独立探测到核球信号，并给出辐射比单一点源更延展的迹象 \cite{Kierans2020,Siegert2020COSI}。谱学测量还约束了湮灭介质的物理状态和运动学：SPI 谱线可分解为约 $1.3\keV$ FWHM 的窄成分和约 $5.4\keV$ 的宽成分，并已在银河系内区搜寻其中心能量与线宽的变化 \cite{Jean2006,Siegert2019Kinematics}。然而，511 keV 天图记录的是正电子减速并湮灭后的位置，不一定是正电子最初产生的位置；正电子在热化之前可在磁化的多相星际介质中传播。因此，现有弥散形态和谱学约束仍不能单独判定主导来源，或还原湮灭前的输运历史 \cite{Prantzos2011}。

宽视场康普顿和编码孔径仪器非常适合测量核球和银盘的总体发射。为了检验中心区域是否还包含未分辨中心增强、紧致源或可测量的谱线运动学结构，需要一种互补的定点仪器。探测到点状或窄集中分布的特征，将表明部分湮灭在空间上是局域的；未探测到信号则会限制紧致成分，但不能排除正电子逃逸后以弥散形式湮灭的源族。INTEGRAL/IBIS 紧致源搜索未发现银河系 511 keV 点源，并针对银河系中心方向约 10 Ms 的曝光给出约 $1.6\times10^{-4}\phcms$ 的 $2\sigma$ 上限 \cite{DeCesare2011}。该上限既给出了以往紧致源搜索的比较尺度，也构成新型窄视场定点仪器的灵敏度目标。

\revmark{1}Laue 光学利用透射式 Bragg 衍射，将来自已知方向的光子聚焦到小型焦面，为定点观测提供有效的集光方式 \cite{Barriere2009Laue,Barriere2011LauePrototype}。当探测器相关本底占主导时，将大口径收集的光子汇聚到较小的灵敏探测器面积可提高点源灵敏度 \cite{Weidenspointner2006MAX}。焦面探测器需要兼具窄能量响应和足够的 511 keV 阻止效率。过渡边缘传感器（transition-edge sensor, TES）微量热计通过超导体陡峭的电阻转变测量能量沉积引起的微小温升，可在 511 keV 附近实现亚 keV 能量分辨率 \cite{Shirazi2023,Zhang2022TESGamma}。阵列用于覆盖焦斑，厚吸收体和多层结构则提高全能峰探测效率。Laue 透镜将目标光子集中到小面积焦面，TES 响应支持窄线窗选取和精细线形测量。由此构成的窄视场聚焦能谱仪以视场范围换取集中的目标光子束和高分辨率谱学能力，可补充以往对核球和银盘的宽视场测量，并用于搜索紧致点源。

能量分辨率和聚焦性能不能单独决定点源灵敏度 \cite{Shirazi2023,Hu2026}。大气辐射场在灵敏体积中的瞬发沉积，以及探测器和低温系统材料活化产生的延迟辐射，也共同决定科学窗内的本底 \cite{Gallego2025Balloon,Tian2026DIXE}。

\revmark{2}本工作主要评估气球载 Laue 透镜--TES 阵列望远镜对银河系中心 511 keV 点源的高空探测性能。我们按物理分量和入射粒子分类科学窗本底，进一步追溯相关材料、母核素和产生区域，以识别主要产生机制，并据此优化探测系统。

\subsection{科学目标与仪器要求}
\label{sec:requirements}

\begingroup
\clubpenalty=10000
\widowpenalty=10000

\revmark{0}本文模拟的仪器用于搜索银河系中心区域的 511 keV 紧致点源和可能尚未分辨的中心超额成分。INTEGRAL/SPI 测量表明，511 keV 天空辐射以弥散的核球和银盘发射为主 \cite{Siegert2016AA,Yoneda2025}。De Cesare 针对 INTEGRAL/IBIS 开展的紧致源专项搜索未发现银河系 511 keV 点源，并给出了银河系中心方向 10 Ms 曝光下的 $2\sigma$ 上限 \cite{DeCesare2011}。为与本文报告的 $3\sigma$ 灵敏度基准比较，我们在高斯近似下对该上限作线性换算，将参考源通量统一定义在大气衰减前：$\fref=2.4\times10^{-4}\phcms$。任务折叠在每个轨迹节点根据假设海拔对应的剩余大气柱深和固定 $45^\circ$ 仰角计算大气透过率。代表性参考状态取第 15 天轨迹节点，海拔为 $38.75\,\mathrm{km}$；该节点采用的剩余大气柱深给出 $T_{\mathrm{ref}}=0.652034$，相应的望远镜入射参考通量为 $F'_{\mathrm{ref}}=\fref T_{\mathrm{ref}}=1.565\times10^{-4}\phcms$。其他轨迹节点处的望远镜入射通量写为 $\fref T_{\mathrm{atm},511}(t)$，最小可分辨通量按大气衰减前的同一口径报告。与 INTEGRAL 的宽视场光谱和编码孔径搜索相比，本文模拟的载荷采用 Laue 聚焦形成点源响应，并以高能量分辨率、多层厚吸收体 TES 微量热计阵列作为焦面探测器 \cite{Shirazi2023,Caputo2017}。

\revdeleted{R4：原“四项要求”段已删除；本句仅为审阅标记。}

本文余下结构如下：第 2 节给出探测器与低温系统质量模型；第 3 节描述模拟框架、事例选择、任务时间轨迹折叠及响应加权环境筛查；第 4 节给出气球环境下的选后率、本底来源和灵敏度；第 5 节给出轨道与月面源场的条件转移；第 6 节讨论物理解释与适用边界，第 7 节总结主要结论。

\endgroup

\section{探测器与低温系统质量模型}
\label{sec:geometry}

本文比较两个完整的探测器--低温系统质量模型：参考方案质量模型 A，以及由质量模型 A 本底分析提出的质量模型 B。两者采用相同的六层 TES 核心，并均包含小型稀释制冷机、混合室、多级冷盘、热屏蔽与磁屏蔽、铜支撑、主动闪烁体和低温服务结构 \cite{Shirazi2023,Gottardi2021TESReview}；但焦面的布置、物理入口、局部 W 结构和主动屏蔽构型并不相同。图~\ref{fig:mass_model} 以相同视角和物理比例比较两种几何。

\revmark{5}TES 探测器概念采用本实验室研发的 AlMn 薄膜与吸收体耦合结构。AlMn 合金的相变温度可通过成分和退火工艺调节，适于阵列制备，并且对磁场相对不敏感 \cite{Xie2026AlMn,Vavagiakis2017AlMnMagnetic}。在本文概念设计中，TES 温度计为厚约 $200\,\mathrm{nm}$、Mn 原子掺杂浓度为 $2000\,\mathrm{ppm}$ 的 AlMn 薄膜，而 $\gamma$ 射线吸收体为毫米量级金属体。吸收体采用 $1.5\times1.5\times3.0\,\mathrm{mm^3}$ Ta 像素；Ta 的高原子序数提高了 511 keV 光子的相互作用效率。在 511 keV 处，单层 $3.0\,\mathrm{mm}$ Ta 的预估相互作用效率约为 $49\%$，相应的六层累计相互作用概率约为 $98\%$。当前几何中，相邻 Ta 阵列层的中心间距为 $12\,\mathrm{mm}$；该间距的设计兼顾后续基于康普顿运动学轨迹重建的 veto 机制。选择 Ta 还因为其在低温区的超导态比热较低。概念热学估算采用 $5.99\times10^{-7}\,\mathrm{J\,kg^{-1}\,K^{-1}}$ 作为代表性低温模型参数。较低的比热可限制毫米级厚吸收体带来的热容增加，使提高阻止效率与保持量热型探测器所需的小热容相兼容，从而有利于能量分辨率。Ta 同时具有较高的光电--康普顿相互作用比；根据 NIST XCOM 在 511 keV 附近的截面，该比值约为 0.731 \cite{NISTXCOM}。在探测器响应模型中，我们根据本实验室已开发的 X 射线 TES 的实验经验 \cite{Xie2026AlMn}，采用 $\Delta E_{\mathrm{FWHM}}\simeq420\,\mathrm{eV}$ 作为 511 keV 处的预估能量分辨率和探测器响应后处理输入。

由于 TES 薄膜与 Ta 吸收体在尺寸和质量上相差很大，且 TES 薄膜本身对 511 keV 光子的相互作用贡献可忽略，输运几何使用 Ta 吸收体代表 TES 像素的主要敏感体积。显式建模的敏感体积共 6 层，每层包含 376 个 Ta 吸收体像素，总计 2256 个像素；每个像素长、宽均为 $1.5\,\mathrm{mm}$，厚度为 $3\,\mathrm{mm}$。每个 Ta 像素下方放置一个 Si 衬底代理；由于本模拟不处理详细热传导，使用 Si 代表 $\mathrm{SiO_2}$--$\mathrm{SiN_x}$--Si 衬底。TES 薄膜、布线和读出细节不作为显式相互作用体积逐层建立，而是通过假设的能量响应和服务质量代理进入模型。与响应 FWHM 独立定义，保留的科学窗选择采用 511 keV 线窗 $W_{511}=511\,\mathrm{keV}\pm420\,\mathrm{eV}$，即 $510.58$--$511.42\,\mathrm{keV}$。

质量模型 A 的 TES 阵列安装在多级冷盘正下方的铜支架上。沿光路依次设置铝热屏蔽窗、Be 真空窗和与 TES 像素对准的 W 多孔准直器；周围壳层通过布尔开孔形成侧入射光路。完整仪器框架旋转 $45^\circ$，使局部孔径在地平坐标中指向 $45^\circ$ 仰角。

\revmark{6}质量模型 A 的来源追溯表明，多级铜冷端结构的活化是科学窗本底的重要来源。该诊断促成了质量模型 B 的完整构型：将同一六层 TES 焦面移入侧向烟囱，并在焦面附近设置局部 BGO。该布局使探测器离开活化冷端结构的直接投影；把 BGO 局限在焦面附近的设计目标，是同时降低耦合本底与主动屏蔽质量，后文评价的是这一完整构型。

\begin{figure}[t]
\centering
\includegraphics[width=0.98\linewidth]{figures/manuscript/fig01_mass_models_ab_zh.pdf}
\caption{质量模型 A（a）与质量模型 B（b）的探测器--低温系统对照图；两面板采用相同的相机、物理视口和比例。质量模型 A 的六层 TES 位于多级冷端结构下方，前方为多孔 W 准直器。质量模型 B 将同一 TES 核心横向移入 Al 烟囱，采用开放四边 W 框，并在烟囱侧面、前端光学环和后端冷孔环设置主动 BGO。红色箭头表示聚焦 511 keV 光子的名义传播路径。为便于观察，部分外层包络设为半透明或隐藏。图中物理入口不等同于拓扑 veto 使用的分析圆盘。}
\label{fig:mass_model}
\end{figure}

对照图直接显示了焦面的横向位移。质量模型 A 中，厚 4.796 mm 的 Bi 上半圆筒安装在 2 mm Al 冷端圆筒内侧，位于六层 Ta 吸收体及其铜热支撑上方；质量模型 B 的六个敏感平面则沿侧烟囱轴布置，离开多级冷盘的直接投影，并位于局部 BGO 组件内。

\revmark{7}两套设计均在同一 MEGAlib/Geant4 模拟与分析框架下计算，采用共同的源定义、探测器响应、科学窗、计时和 Compton 选择，并按各自几何施加主动反符合；所得性能比较对应两套完整构型。

\FloatBarrier

\section{模拟框架}
\label{sec:framework}

\revmark{8}本文采用 MEGAlib（Medium Energy Gamma-ray Astronomy Library）开展模拟；该框架以 Geant4 为基础，用于中能 $\gamma$ 射线仪器的建模与事例分析 \cite{Zoglauer2006,Agostinelli2003,Allison2016}。其中 Cosima 模块定义粒子源和探测器几何，完成质量模型中的粒子输运，并记录后续探测器响应与事例选择所需的相互作用信息。图~\ref{fig:workflow} 将端到端计算按物理来源分为三条分支：聚焦信号、瞬发大气本底和延迟活化本底。Laue 光学系统与探测器--低温系统分别建立质量模型。参考点源首先在 Laue 透镜模型中输运；到达焦平面的光子位置与方向写入 MEGAlib EventList，并据此初始化探测器侧 Monte Carlo 输运。对于每一版设计，聚焦信号输入与全部本底源均在第~\ref{sec:geometry}~节所述的同一完整探测器--低温系统质量模型中输运。

瞬发输运与活化产生输运并行采用同一套\revblue{ EXPACS/PARMA 八族宽带辐射场，包含 $\gamma$、$e^-$、$e^+$、$\mu^-$、$\mu^+$、$n$、$p$ 和 $\alpha$}。在延迟活化分支中，BUILDUP 输运记录核素产额、材料体积和产生位置；这些记录用于计算飞行期间积累的活度，并在各参考时刻构建核素清单。随后按记录的产生位置抽样延迟衰变，并在探测器--低温系统模型中输运。

聚焦信号与当前率预算中的\revblue{两类本底流——八类瞬发大气粒子族和延迟活化——}在每个质量模型内统一施加探测器响应。\revmark{9}响应后，全部信号与本底流采用 $510.58\leq E_{\mathrm{TES}}<511.42\keV$ 的科学窗；窗内候选随后接受各自几何的主动 veto 和共同 Compton 拓扑检验。最终保留的本底事例按入射粒子族、母核素、材料和产生体积分解。来源诊断用于分析质量模型 A，并据此提出和检验完整备选构型 B。两套设计固定后，其信号率与分量本底率沿共同轨迹折叠，得到累计计数和未分辨线灵敏度。下文给出源模型、输运归一化、事例选择和任务时间计算。

\begin{landscape}
\centering
\vspace*{\fill}
\includegraphics[width=0.94\linewidth]{figures/manuscript/fig02_workflow_zh_threebranch_20260830.pdf}
\captionof{figure}{\revmark{17}端到端模拟与分析流程。聚焦信号、瞬发大气本底和延迟活化构成三条物理分支；\revblue{瞬发分支包含八个独立归一的粒子族（$\gamma,e^-,e^+,\mu^-,\mu^+,n,p,\alpha$）}。每个质量模型内的全部事例流统一施加探测器响应、主动 veto 和 Compton 拓扑检验。选后事例来源诊断用于提出完整的质量模型 B 构型，两套固定设计随后在共同任务时间轴上评价。}
\label{fig:workflow}
\vspace*{\fill}
\end{landscape}

\FloatBarrier

\subsection{Laue 光学与探测器耦合信号输运}
\label{sec:optics}

本文模拟采用镶嵌型 Ge(111) 晶体 Laue 透镜 \cite{Barriere2009Laue,Barriere2011LauePrototype}。该光学设计的焦距为 $10\,\mathrm{m}$，在 511 keV 处形成单环参考响应，有效面积为 $20.1\,\mathrm{cm^2}$，并按
\begin{equation}
  A_{\mathrm{eff}}(E)=A_{\mathrm{inc}}\frac{N_{\mathrm{fp}}(E)}{N_{\mathrm{inc}}(E)}
  \label{eq:laue_aeff}
\end{equation}
计算，其中 $A_{\mathrm{inc}}$ 为蒙特卡洛入射采样面积，$N_{\mathrm{inc}}(E)$ 为入射光子数，$N_{\mathrm{fp}}(E)$ 为满足聚焦/焦面截取条件的光子数。
镶嵌晶体给出设计所需的角接受度，并在全部光学输运中采用同一实现。

Laue 光学分支同时模拟聚焦过程及光子在 Ge 晶体中的输运。

为此，晶体衍射概率不作为额外效率因子叠加，而在 Geant4 应用层定义为离散过程 \cite{Agostinelli2003,Allison2016}。该过程根据光子相对晶面的角偏差取得衍射概率 $R(|\Delta\theta|)$，再将其转换为有限的等效衍射平均自由程 $\lambda_{\mathrm{diff}}$，使光子穿过晶体路径 $\ell$ 时的累计衍射概率满足 $1-\exp(-\ell/\lambda_{\mathrm{diff}})=R(|\Delta\theta|)$，即 $\lambda_{\mathrm{diff}}=-\,\ell/\ln[1-R(|\Delta\theta|)]$。为避免与标准光电吸收重复计数，$R$ 先按未吸收份额归一为 $R/(1-A)$，再转换为平均自由程。Geant4 随后在同一竞争框架下分别抽样衍射、光电吸收和康普顿散射的相互作用长度，并由最短者决定首先发生的过程，而不是在晶体边界强制截获光子。未发生衍射的光子继续按标准电磁物理输运，可能被吸收、散射或直接透射。

本文采用预先计算的 XOP/CRYSTAL rocking curve 确定衍射概率。Ge(111) 的反射率、透射率和吸收率随角偏差的关系由 XOP 工具包的 CRYSTAL/diff\_pat 模块生成（通过 xoppylib 调用，晶体学参数取自 DABAX）\cite{SanchezDelRio2011XOP}，并在预处理阶段存为输运查找表。在 Geant4 输运过程中，对晶体内的每个输运步计算光子相对于 Bragg 条件的角偏差，并通过插值得到相应的衍射概率。30 角秒 FWHM 的镶嵌角分布通过对理想 Bragg 晶面法线施加高斯角扰动实现；发生衍射后，根据扰动后的晶面法线确定出射方向，光子能量保持不变。该方法借鉴了已有的 Geant4 晶体 Bragg 反射实现 \cite{Guan2023BraggReflection,Reiazi2025G4BraggReflection}，并应用于本文研究的能段。

这里的参考信号是单色 511 keV 点源，采用定义在大气层外的唯一参考通量 $\fref=2.4\times10^{-4}\phcms$（见第~\ref{sec:requirements} 节）。该尺度来自 INTEGRAL/IBIS 的紧致源专项上限，而非某个已观测天体 \cite{DeCesare2011}。信号以 MEGAlib 远场源方式构建 \cite{Zoglauer2006}，详细构建方式见第~\ref{subsec:pointsource} 节。信号经光学分支输运并作为平行束源输入探测器质量模型；宇宙线瞬发本底与活化衰变则直接在探测器/低温系统质量模型中输运。因此，主要本底预算定义在探测器/低温系统分支上，光学分支提供输运到焦平面的信号相空间。

\subsection{源构建}
\label{sec:background}

参考点源的焦面 EventList、瞬发大气场和由活化核素清单构建的延迟衰变源彼此独立。\revblue{瞬发场由八个粒子族分别归一}，活化产生与延迟衰变则作为两个独立输运阶段处理。该组织方式在施加共同探测器响应和事例选择前，保留了各本底流的归一化与来源追溯关系。

\subsubsection{瞬发大气源}
\label{subsec:prompt_source}

瞬发大气粒子场由 EXPACS/PARMA 给出。PARMA3.0 是由 PHITS 大气级联计算参数化得到
并针对陆地宇宙线测量验证过的解析模型，EXPACS 是其公开表格与软件接口
\cite{Sato2015EXPACS,Sato2016PARMA4}。底层的粒子输运物理由 PHITS 提供：一次宇宙线
与大气核作用，次级强子、轻子、光子和中子在大气柱中继续输运，最终通量按海拔、
地磁截止刚度、太阳调制、能量和方向参数化。本文中 EXPACS/PARMA 给出随粒子种类、
动能、天顶角、地理位置、海拔、截止刚度和太阳活动变化的入射微分通量，且只用于构造
外部辐射场；探测器响应、反符合、拓扑选择和任务时间计数都在下游施加。参考源定义
采用纬度 $34^\circ$、经度 $100^\circ$、海拔 $38\,\mathrm{km}$、垂直截止刚度
$R_c=11.6\,\mathrm{GV}$，以及 2025-08-31 对应的太阳条件 $W=118.3$。

瞬发源包含光子、中子、电子、正电子、质子、$\alpha$ 粒子以及正负 $\mu$ 子。对每个
粒子族，EXPACS/PARMA 微分谱被转换为 MEGAlib 远场面积源，覆盖完整天顶角范围
$\theta=0^\circ$--$180^\circ$ 和完整方位角范围 $\phi=0^\circ$--$360^\circ$。天顶角
依赖用 20 个等 $\mu$ 微分分箱表示，其中 $\mu=\cos\theta$，每箱立体角
$\Delta\Omega=0.6283185307\,\mathrm{sr}$：前 10 箱为下行粒子，后 10 箱为上行或反照
成分。瞬发输运和活化产生输运都保留这两个半球，因此源构建阶段没有丢弃上行反照
粒子。图~\ref{fig:expacs_fullsphere_flux} 给出由源定义实际采用的全空间通量场。

\begin{figure}[t]
\centering
\includegraphics[width=\linewidth]{figures/manuscript/fig03_corrected_kev_sources.pdf}
\caption{用于瞬发输运和活化产生计算的 EXPACS/PARMA 大气粒子场。（a）八类粒子族在下行和上行（反照）半球内对能量和立体角积分后的通量。（b）各粒子族在 20 个等 $\mu$ 天顶角分箱（$\mu=\cos\theta$）中的积分通量。环境参数为纬度 $34^\circ$、经度 $100^\circ$、海拔 $38\,\mathrm{km}$、$R_c=11.6\,\mathrm{GV}$ 和 $W=118.3$。}
\label{fig:expacs_fullsphere_flux}
\end{figure}

全空间积分通量分别为 4.79966、0.462239、0.199002、0.116981、0.112301、0.0114871、0.00496893 和 $0.00557001\,\mathrm{cm^{-2}\,s^{-1}}$，依次对应 $\gamma$、$n$、$e^-$、$e^+$、$p$、$\alpha$、$\mu^-$ 和 $\mu^+$。每一族均采用 20 个等 $\mu$ 分箱和半径 60 cm 的远场起始球。

\revdeleted{R10：原源能量单位换算说明已删除；本句仅为审阅标记。}

若 $I_j(E,\mu)$ 表示粒子族 $j$ 的 EXPACS/PARMA 微分强度，则等 $\mu$ 分箱 $m$ 承载的通量为
\begin{equation}
 \Phi_{j,m}=2\pi\int_{\mu_m}^{\mu_{m+1}}\!\mathrm d\mu
             \int I_j(E,\mu)\,\mathrm dE,
 \qquad \Delta\mu=0.1,
 \label{eq:source_bin_flux_zh}
\end{equation}
故每箱立体角为 $\Delta\Omega=2\pi\Delta\mu=0.62831853\,\mathrm{sr}$。前 10 箱覆盖下行半球，后 10 箱覆盖上行或反照半球，两者方位角均完整。

八类粒子族均采用半径 $R=60\,\mathrm{cm}$ 的远场面积源。对每个抽样方向，起点在垂直于该方向的投影圆盘内均匀抽样，动量指向质量模型内部，因而几何率因子为投影面积 $\pi R^2$。20 个角分箱共同覆盖 $4\pi$ 入射立体角。同一大气粒子场定义分别用于瞬发输运和活化产生计算。

\begingroup
\color{revisionblue}
\noindent\textbf{[R14]} 每个粒子族 $j$ 的已接受生产任务都同时包含完整的 20 个角分箱混合，因而这些角分箱不作为彼此独立的曝光累加。该族的源端初级粒子率和合并等效照射时间定义为
\begin{equation}
 \dot N^{\mathrm{src}}_{j}=\pi R^2\sum_{m=1}^{20}\Phi_{j,m},
 \qquad
 T^{\mathrm{eq}}_{g,j}=\sum_{q\in\mathcal Q_{g,j}}\mathrm{TT}_{g,j,q}
 \simeq
 \frac{\sum_q N^{\mathrm{sim}}_{g,j,q}}{\dot N^{\mathrm{src}}_{j}},
 \qquad
 w^{\mathrm{pr}}_{g,j}=\frac{1}{T^{\mathrm{eq}}_{g,j}},
 \label{eq:prompt_equivalent_exposure_zh}
\end{equation}
其中 $q$ 标记彼此独立且可合并的任务，$\mathrm{TT}_{g,j,q}$ 是相应 receipt 中校验通过的物理运行时间。增加输运初级粒子数会延长 $T^{\mathrm{eq}}$ 并降低 Monte Carlo 统计误差，而不改变物理源强。七类非 $\gamma$ 瞬发粒子的探测器正沉积模板继承常数权重 $w^{\mathrm{pr}}_{g,j}$；宽带 $\gamma$ 目录则保留逐事例谱形重要性因子。送入公共时间轴泊松抽样的探测器正沉积目录率为
\begin{equation}
 R^{\mathrm{cat}}_{g,j}=\sum_{e\in\mathcal C^{\mathrm{det+}}_{g,j}}w_e,
 \qquad
 N_{g,j}(\Delta t)\sim
 \operatorname{Poisson}\!\left(R^{\mathrm{cat}}_{g,j}\Delta t\right),
 \label{eq:catalog_poisson_rate_zh}
\end{equation}
其中 $\mathcal C^{\mathrm{det+}}_{g,j}$ 包含在 TES 或相应几何的主动屏蔽体积中产生正原始能量沉积的事例。该占用率定义在 420 eV FWHM 响应、0.3 keV 像素阈值、科学窗、主动 veto 和 Compton 拓扑选择之前；谱形重要性因子在求和前已包含于 $w_e$。
\endgroup

\subsubsection{延迟活化源}
\label{subsec:delayed_source}

延迟源来自独立的活化产生输运。活化记录不作为即时本底计数，而是保存每个放射性产物的核素、材料体积、入射粒子族和产生位置；基态半衰期对照 NUBASE2020 校验 \cite{Kondev2021NUBASE}。对核素 $k$，
\begin{equation}
  A_k(t_{\mathrm{flight}})=P_k\left[1-\exp(-\lambda_k t_{\mathrm{flight}})\right],
  \qquad \lambda_k=\frac{\ln2}{t_{1/2,k}},
\end{equation}
其中 $P_k$ 为产生率，活化模拟以气球飞行第 15 天为参考状态，即
$t_{\mathrm{flight}}=15\,\mathrm{d}$。

\begingroup
\color{revisionblue}
\noindent\textbf{[R11]} 在本文采用的标准输出配置下，MEGAlib/Cosima 的压缩 \texttt{sim.gz} 事例文件未保留重建延迟源空间分布所需的完整放射性产物产生信息。Cosima 以 Geant4 为输运引擎，因此本文在 \texttt{SteppingAction} 中加入记录钩子，在放射性产物生成时保存其产生坐标、材料体积、核素和激发态。随后利用这些字段将第 15 天核素清单与产生记录匹配，并按入射粒子族、材料类别、产生体积和坐标组织延迟源。

保留核素产生位置十分重要，因为延迟粒子能否穿出周围材料、触发相应几何的主动反符合或通过 Compton 拓扑选择，均取决于产生地点。衰变位置按活度从匹配的产生记录中抽样，使延迟目录同时保持空间分布和总活度。两种质量模型的活化源位置及其中影响 TES 511 keV 响应窗的母核素示于图~\ref{fig:activation_source_positions_zh}。
\endgroup

\revdeleted{R12：原有限选后事例数的重复陈述已删除；本句仅为审阅标记。}

\begingroup
\color{revisionblue}
\noindent\textbf{[R15]} 每个质量模型对每一入射粒子族输运 $N^{\mathrm{dec}}_{g,j}=10^6$ 个延迟衰变。以第 15 天为参考时刻，该族延迟源的恒定活度等效照射时间和单事例率权重为
\begin{equation}
 T^{\mathrm{eq,del}}_{g,j}(t_{\mathrm{ref}})
 =\frac{N^{\mathrm{dec}}_{g,j}}{A_{g,j}(t_{\mathrm{ref}})},
 \qquad
 w^{\mathrm{del}}_{g,j}(t_{\mathrm{ref}})
 =\frac{A_{g,j}(t_{\mathrm{ref}})}{N^{\mathrm{dec}}_{g,j}}
 =\frac{1}{T^{\mathrm{eq,del}}_{g,j}(t_{\mathrm{ref}})},
 \quad t_{\mathrm{ref}}=15\,\mathrm{d}.
 \label{eq:delayed_equivalent_exposure_zh}
\end{equation}
这里的 $T^{\mathrm{eq,del}}$ 是在第 15 天活度混合下的 Monte Carlo 等效抽样时间，而不是实际飞行时长或 BUILDUP 照射时间。探测器侧延迟目录率仍按式~(\ref{eq:catalog_poisson_rate_zh})对通过输运的模板权重求和；任务折叠则按母核素活度曲线将参考权重更新到各时间节点。
\endgroup

\begingroup
\color{revisionblue}
\noindent\textbf{[R13]} 对于 511 keV 附近的活化本底，两套模型的最终事例均主要关联制冷系统中的铜结构。质量模型 A 的来源分解表明，TES 邻近冷盘和混合室附近的铜是重要贡献区。基于这一诊断，质量模型 B 将焦面布置在侧向烟囱内，以减弱 TES 与主要铜活化区的几何耦合，并在冷端方向重新布置局部 BGO，从而提高相关事例产生符合反符合标记的概率。在共同源场、探测器响应和选择条件下，完整质量模型 B 的延迟本底低于质量模型 A。由于入口、邻近材料和主动屏蔽布局同时变化，现有比较不能把该差异单独归因于空间位移或 BGO 布局中的任一项；第~\ref{sec:results} 节将进一步讨论来源与几何的关系。
\endgroup

\begingroup
\color{revisionblue}
\begin{table}[!htbp]
\color{revisionblue}
\centering
\captionsetup{font={color=revisionblue}}
\caption{\textbf{[R16]} 两种质量模型的粒子族输运统计与率归一。$N_{\mathrm{tr}}$ 和 $N_{\mathrm{dec}}$ 分别为输运的瞬发初级粒子数和延迟衰变数，$T_{\mathrm{eq}}$ 为已接受独立生产批次的物理曝光之和，$N_{\mathrm{det+}}$ 为在 TES 或主动屏蔽中具有正原始能量沉积的模板数。$R_{\mathrm{cat}}$ 是实际送入公共时间轴泊松抽样的探测器正沉积占用率；其定义位于探测器响应展宽、像素阈值、科学窗、主动 veto 和 Compton 选择之前。面板（b）采用 $T_{\mathrm{eq,del}}=N_{\mathrm{dec}}/A_{g,j}(15\,\mathrm d)$，定义见式~(\ref{eq:delayed_equivalent_exposure_zh})。}
\label{tab:background_source_model}
\scriptsize
\setlength{\tabcolsep}{5.0pt}
\renewcommand{\arraystretch}{0.92}
\noindent\textit{(a) 瞬发大气粒子输运}\par\vspace{0.20em}
\begin{tabular}{clrrrr}
\toprule
模型 & 粒子族 & $N_{\mathrm{tr}}$ & $T_{\mathrm{eq}}$ (s)
& $N_{\mathrm{det+}}$ & $R_{\mathrm{cat}}$ ($\mathrm{s^{-1}}$) \\
\midrule
A & $\gamma$ & $20{,}315{,}706$ & 374.249513 & $6{,}372{,}612$ & $1.635799\times10^{4}$ \\
A & $n$ & $1{,}164{,}960$ & 223.112900 & $403{,}599$ & $1.808945\times10^{3}$ \\
A & $e^-$ & $675{,}057$ & 299.757900 & $287{,}509$ & 959.137357 \\
A & $e^+$ & $394{,}119$ & 298.875770 & $182{,}318$ & 610.012648 \\
A & $p$ & $380{,}829$ & 299.752500 & $172{,}734$ & 576.255411 \\
A & $\alpha$ & $39{,}050$ & 298.957500 & $18{,}657$ & 62.406864 \\
A & $\mu^-$ & $16{,}607$ & 298.151700 & $7{,}529$ & 25.252246 \\
A & $\mu^+$ & $19{,}371$ & 308.674500 & $8{,}629$ & 27.955014 \\
\addlinespace[0.25em]
B & $\gamma$ & $112{,}421{,}262$ & 2071.102900 & $4{,}189{,}660$ & $1.939983\times10^{3}$ \\
B & $n$ & $8{,}363{,}786$ & 1599.414700 & $325{,}382$ & 203.438170 \\
B & $e^-$ & $3{,}602{,}914$ & 1598.881500 & $136{,}738$ & 85.521035 \\
B & $e^+$ & $2{,}119{,}548$ & 1600.953780 & $139{,}963$ & 87.424760 \\
B & $p$ & $2{,}036{,}710$ & 1602.777690 & $123{,}624$ & 77.131096 \\
B & $\alpha$ & $209{,}074$ & 1607.562620 & $18{,}285$ & 11.374363 \\
B & $\mu^-$ & $91{,}020$ & 1608.091082 & $4{,}605$ & 2.863644 \\
B & $\mu^+$ & $102{,}306$ & 1612.825350 & $5{,}268$ & 3.266318 \\
\bottomrule
\end{tabular}

\vspace{0.45em}
\noindent\textit{(b) 第 15 天延迟衰变输运}\par\vspace{0.20em}
\begin{tabular}{clrrrr}
\toprule
模型 & 粒子族 & $N_{\mathrm{dec}}$ & $T_{\mathrm{eq,del}}$ (s)
& $N_{\mathrm{det+}}$ & $R_{\mathrm{cat},15}$ ($\mathrm{s^{-1}}$) \\
\midrule
A & $\gamma$ & $1.0\times10^6$ & $1.516919\times10^5$ & $884{,}678$ & 5.832072 \\
A & $n$ & $1.0\times10^6$ & 3030.110 & $893{,}664$ & 294.927907 \\
A & $e^-$ & $1.0\times10^6$ & $1.497281\times10^6$ & $853{,}592$ & 0.570095 \\
A & $e^+$ & $1.0\times10^6$ & $3.385589\times10^5$ & $904{,}775$ & 2.672430 \\
A & $p$ & $1.0\times10^6$ & 1380.694 & $916{,}824$ & 664.031386 \\
A & $\alpha$ & $1.0\times10^6$ & 3477.763 & $918{,}773$ & 264.185032 \\
A & $\mu^-$ & $1.0\times10^6$ & $4.034086\times10^6$ & $871{,}246$ & 0.215971 \\
A & $\mu^+$ & $1.0\times10^6$ & $7.933795\times10^{14}$ & $1{,}000{,}000$ & $1.260431\times10^{-9}$ \\
\addlinespace[0.25em]
B & $\gamma$ & $1.0\times10^6$ & $4.921159\times10^5$ & $396{,}106$ & 0.804904 \\
B & $n$ & $1.0\times10^6$ & $1.459182\times10^4$ & $554{,}831$ & 38.023436 \\
B & $e^-$ & $1.0\times10^6$ & $3.754719\times10^6$ & $433{,}022$ & 0.115327 \\
B & $e^+$ & $1.0\times10^6$ & $1.480682\times10^6$ & $531{,}844$ & 0.359189 \\
B & $p$ & $1.0\times10^6$ & 9476.321 & $694{,}056$ & 73.241080 \\
B & $\alpha$ & $1.0\times10^6$ & $2.916849\times10^4$ & $711{,}199$ & 24.382444 \\
B & $\mu^-$ & $1.0\times10^6$ & $1.461788\times10^7$ & $436{,}535$ & 0.029863 \\
B & $\mu^+$ & $1.0\times10^6$ & $2.521327\times10^8$ & $1{,}867$ & $7.404830\times10^{-6}$ \\
\bottomrule
\end{tabular}

\vspace{0.35em}
\parbox{0.97\linewidth}{\raggedright 宽带 $\gamma$ 瞬发目录的 $R_{\mathrm{cat}}=\sum_e w_e$ 包含逐事例谱形重要性权重；其余七类瞬发粒子采用常数 $1/T_{\mathrm{eq}}$ 权重。延迟权重在第 15 天同一族内为常数，任务折叠则逐母核素演化活度，而不是用单一因子整体缩放一个粒子族。}
\end{table}
\endgroup

\begin{figure}[p]
\centering
\includegraphics[width=0.96\linewidth]{figures/manuscript/fig_activation_positions_model_a_zh.pdf}
\vspace{0.35em}

\includegraphics[width=0.96\linewidth]{figures/manuscript/fig_activation_positions_model_b_zh.pdf}
\caption{质量模型 A（上）和质量模型 B（下）的延迟活化产生位置。每行左图为质量模型全景，右图为 TES 区域放大；蓝色点表示全部输运的延迟活化源位置，彩色点表示在 TES $510.58$--$511.42$ keV 响应窗中产生本底的活化母核素。}
\label{fig:activation_source_positions_zh}
\end{figure}

\FloatBarrier

\subsubsection{参考点源构建}
\label{subsec:pointsource}

\begingroup
\color{revisionblue}
\revmark{18}聚焦信号采用轴上单色 $511\keV$ 远场点源，在透镜入射面表现为平行光束。参考源的通量定义和大气衰减约定见第~\ref{sec:requirements} 节；本节只说明焦面光子相空间的构建及探测器有效面积的计算。

在光学输运中，入射光子在式~\eqref{eq:laue_aeff} 定义的透镜孔径 $A_{\mathrm{inc}}$ 上均匀抽样。焦面响应由 30 角秒晶体 mosaic 展宽、XOP/CRYSTAL 摇摆曲线和 10 m 聚焦几何共同确定。三个独立光学样本共输运 $1.5\times10^5$ 个入射光子，得到 37,256 个焦面穿越，其中 $N_{\mathrm{fp}}=37{,}194$ 个落入半径为 $1.898\,\mathrm{cm}$ 的 Be 入射窗。相应光斑的包围半径为 $r_{90}=1.03\,\mathrm{cm}$ 和 $r_{99}=1.25\,\mathrm{cm}$。每个被入口接受的光子均以其焦面位置和传播方向写入共同的 EventList，并分别在质量模型 A 和质量模型 B 中独立输运。

该焦面 EventList 表示的光学有效面积为 $A_{\mathrm{fp}}\equiv A_{\mathrm{eff}}(511\keV)=20.08476\,\mathrm{cm^2}$，因此每个等权焦面光子的面积权重为
\begin{equation}
 a_{\mathrm{ray}}=\frac{A_{\mathrm{fp}}}{N_{\mathrm{fp}}}.
 \label{eq:focal_eventlist_area_weight_zh}
\end{equation}
对于质量模型 $g$，设 $N_{\mathrm{sel},g}$ 为经过探测器输运、响应卷积和科学窗筛选，并通过相应构型的孤立事例主动反符合判据以及 Compton 拓扑/孔径选择后保留的焦面光子数。两模型采用共同的 Compton 重建算法，但孔径相容性根据各自的物理入口判断。轴上孤立事例的最终有效面积为
\begin{equation}
 A_{\mathrm{eff},g}
 =\sum_{i\in\mathrm{sel},g}a_{\mathrm{ray}}
 =A_{\mathrm{fp}}\frac{N_{\mathrm{sel},g}}{N_{\mathrm{fp}}}.
 \label{eq:selected_signal_effective_area_zh}
\end{equation}
这一有效面积不包含聚焦信号与独立本底事例偶然重叠造成的附加损失。已经通过孤立事例选择的弱信号探针在公共本底时间轴上的条件存活率由五个时间锚点单独计算，并在任务时间折叠中施加。各轨迹节点的物理信号率由这里得到的有效面积和第~\ref{sec:requirements} 节定义的载荷处入射通量共同确定。

上述有效面积只描述轴上点源响应。离轴角和飞行姿态抖动会通过 Laue 透镜有限的角接受度改变聚焦效率和焦斑位置，本文将其作为信号有效面积的指向系统学边界 \cite{Barriere2009Laue,Weidenspointner2006MAX}。当前任务计算采用固定 $45^\circ$ 仰角的设计光轴，焦面 EventList 沿该轴从各模型的物理入口注入；该轴上有效面积不直接代表扩展源或任意离轴方向的响应。
\endgroup

\revdeleted{R18：删除了本小节中重复的参考通量、大气透过率和第 15 天性能结果；通量折叠保留在任务时间计算中，A/B 数值保留在第 4 节结果中。}

\subsection{探测器响应、事例选择与 veto 定义}
\label{sec:selection}

公共时间计算先把瞬发大气本底与延迟活化本底目录映射到同一条显式时间轴，再顺次施加主动 veto 和 Compton 选择。\revblue{瞬发流内的八类粒子族各自保留物理率归一化}。聚焦信号作为条件探针逐个插入，用于计算偶然符合损失，而不计入本底预算。

\subsubsection{时间轴与率归一}
\label{subsec:time_axis_zh}

在固定辐射环境下，如果粒子事例彼此独立且平均到达率保持稳定，则给定时间窗内的事例数服从泊松分布；在事例数给定后，各事例的到达时刻可视为在该时间窗内均匀分布。这为率归一的 Monte Carlo 事例目录提供了自然的时间表示：事例目录给出能量沉积、空间位置和拓扑等探测器信息，泊松过程则决定这些事例在物理时间轴上出现的次数和时刻。

气球飞行中的辐射场会随轨迹变化，但这种变化远慢于本文采用的 $1\,\mu\mathrm{s}$ 符合时间尺度。因此，在每个轨迹节点内，各类事例流均可近似为到达率恒定且相互独立的泊松过程。聚焦信号、瞬发大气本底和延迟活化本底由不同的输运计算产生，原始目录之间不共享物理时间轴；而相应几何的主动反符合和康普顿拓扑选择取决于同一时间窗内的全部能量沉积。为计算不同来源之间的偶然符合，需要先将这些目录映射到同一条时间轴。

三个输运事例流在第~\ref{sec:background} 节的源归一化阶段各自获得逐记录率权重 $r_{km}$。其中 $k$ 取聚焦信号、瞬发大气本底和延迟活化本底三类，$m$ 标记该目录中的输运记录，目录总率为 $R_k=\sum_m r_{km}$。对于长度为 $T$ 的参考观测窗，第 $k$ 个目录产生的事例数由下式抽取：
\begin{equation}
 N_k\sim\mathrm{Pois}(R_kT).
 \label{eq:poisson_superposition_zh}
\end{equation}
随后以 $r_{km}/R_k$ 为概率有放回地抽取记录，并在 $[0,T]$ 内为每个实例赋予独立均匀时刻。这样生成的时间序列同时保持了事例流的物理率及其能量、方向和探测器响应分布。

三组实例合并后按到达时刻排序。相邻事例的时间间隔不超过 $\tau=1\,\mu\mathrm{s}$ 时，将其递归归入同一个符合候选组 $c$；组内的全部能量沉积共同接受科学窗、相应几何的主动反符合和康普顿拓扑选择。图~\ref{fig:poisson_time_axis_schematic} 给出这一过程的示意。

\begin{figure}[H]
\centering
\includegraphics[width=\linewidth]{figures/manuscript/fig04a_poisson_time_axis_schematic_threebranch_20260830.pdf}
\caption{公共符合时间轴的构建示意。（a）每个输运事例流具有逐记录率权重和总率；（b）独立泊松抽样合并并排序到公共时间轴；（c）间隔不超过 $\tau=1\,\mu\mathrm{s}$ 的事例组成一个符合候选组。颜色表示目录身份，事例位置和密度均为示意。瞬发与延迟流构成本底时间轴，聚焦信号作为条件探针插入。}
\label{fig:poisson_time_axis_schematic}
\end{figure}

上述计算在第 0、5、10、15 和 20 天五个参考时间点分别进行。在每个参考时间点，根据\revblue{八类瞬发粒子族和延迟活化}各自的物理率进行泊松抽样，组合出该参考环境下的公共本底时间轴。时间间隔不超过 $1\,\mu\mathrm{s}$ 的事例归入同一个符合组，并共同接受科学窗、相应几何的主动反符合和康普顿拓扑选择。聚焦信号事例也逐个加入相应的本底时间轴，用于计算偶然符合造成的信号损失。

随后，本文沿气球的长期飞行轨迹计算仪器性能，并将 20 d 轨迹按 6 h 间隔离散。粒子通量、信号大气透过率和活化库存均随飞行时间及空间位置变化，这些变化也会改变公共时间轴上的事例密度和偶然符合概率。五个参考时间点得到的本底修正和信号存活率用于描述这种变化，并应用于完整轨迹上的信号和本底计算。具体的轨迹模型、时间插值和累计计数方法见第~\ref{sec:mission_time_fold} 节。

\subsubsection{探测器响应、科学窗与带权率}
\label{subsec:detector_response_zh}

在进行科学窗和反符合判断前，同一候选组内的 TES 能量沉积先按像素合并。同一像素中的多次沉积相加为该像素的真实沉积能量，并以能量加权位置代表其作用位置。本文采用预估的单像素能量分辨率 $\Delta E_{\mathrm{FWHM}}=420\,\mathrm{eV}$，并以具有该 FWHM 的高斯响应近似 TES 的测量涨落。展宽后测量能量低于 $0.300\keV$ 的像素命中被移除。质量模型 A 和质量模型 B 采用相同的响应与阈值设置。

所有保留 TES 像素命中的测量能量相加，得到事例的 TES 总测量能量，用于 511\,keV 科学窗判断；各像素的能量、位置、层号和多重度仍分别保留，用于后续康普顿拓扑选择。BGO 中的沉积能量独立累计，不计入 TES 总能量，并在下一小节用于主动反符合判断。

当前主分析统一采用 $510.58\le E_{\mathrm{TES}}<511.42\keV$ 的科学窗。该窗口是根据预估能量分辨率预先固定的高包含率基准窗，对聚焦信号和所有本底分量采用相同定义，并在两个质量模型之间保持不变，避免根据不同模型的结果分别调整选择条件。

本文给出的流量阈值适用于本征宽度相对探测器能量分辨率较小、经响应卷积后仍可由该科学窗充分包含的谱线；若源线明显展宽，须将源线型与响应卷积并重新优化窗宽。

不同源目录具有不同的物理通量、输运曝光和源级重权，因此原始记录数只反映 Monte Carlo 输运支撑，不能直接代表物理率。每条事例均携带相应的物理率权重，各选择阶段的信号率和本底率由通过该阶段的事例权重累加得到。有限输运统计误差按实际权重分布传播，聚焦信号接受度按共同焦面样本的二项统计处理。任务轨迹上的长期率折叠见第~\ref{sec:mission_time_fold} 节；经过主动反符合与拓扑/孔径两级选择的响应卷积能谱和 cut-flow 结果一并在第~\ref{subsec:optimized_fullchain_zh} 节给出。

\subsubsection{主动反符合门}

主动反符合用于排除在 TES 中形成 511\,keV 候选能量、同时在外围主动 BGO 中留下能量沉积的本底事例。被动屏蔽通过材料衰减入射辐射，主动 BGO 则为逐事例判别提供时间符合信息。聚焦 511\,keV 光子主要沿光学通道到达 TES；大气带电粒子、强子相互作用产生的次级粒子以及部分宽带 $\gamma$ 射线级联更容易同时在 TES 和外围主动体中产生沉积。因此，BGO 中的符合沉积可用于识别科学窗内的本底候选。

主动反符合沿用第~\ref{subsec:time_axis_zh} 节定义的公共时间轴和 $1\,\mu\mathrm{s}$ 符合分组。对于每个候选组，TES 沉积按照第~\ref{subsec:detector_response_zh} 节的方法先按像素合并，再施加能量响应，并由所有保留像素的测量能量之和判断是否落入科学窗。同组内全部主动 BGO 体积中的模拟能量沉积相加，记为 $E_{\mathrm{BGO}}$。该沉积只用于反符合判断，不计入 TES 事例能量。

质量模型 A 的主动系统包括 BGO 侧壳、底盖和顶部开孔环；质量模型 B 的主动系统由烟囱侧屏、前端光学环和后端冷孔环三个 BGO 体积组成。两套质量模型采用相同的组级通过条件：
\[
 E_{\mathrm{BGO}}<50\keV .
\]
阈值作用于整个符合组内的 BGO 沉积总和；沉积达到或超过 $50\keV$ 时即否决该候选组。对于给定质量模型，聚焦信号、\revblue{八类瞬发粒子族和延迟活化}均采用相同的组级判据。

组级求和同时包含同一入射粒子在 TES 与主动屏蔽中产生的相关沉积，以及相互独立的命中因满足 $1\,\mu\mathrm{s}$ 链接条件而形成的偶然重叠。聚焦光子若与无关的主动体命中进入同一符合组，也可能被反符合门拒绝；相应的信号损失由第~\ref{subsec:time_axis_zh} 节定义的偶然符合存活率表征，并在任务时间折叠中计入。通过主动反符合的候选事例随后进入康普顿拓扑和入射孔径相容性选择。

本分析直接使用 Cosima 输运记录中的模拟能量沉积，并将 $50\keV$ 作为统一的离线反符合阈值。

\subsubsection{康普顿拓扑与入射孔径相容性}

主动反符合首先排除在外围闪烁体中沉积能量达到所用阈值的候选事例，却无法区分聚焦光子和完全在 TES 阵列内发生相互作用的本底；二者都可能在焦面留下接近 511\,keV 的总能量。不过，聚焦光子应从侧向光学通道进入探测器舱。对于包含两个或更多空间分离命中的事例，各命中的能量分配和相对位置能够约束可能的入射方向。本文只用这些信息回答一个问题：该事例是否可能来自入射孔径？这一判据沿用康普顿相机的运动学成锥方法~\cite{Zoglauer2006}，但这里只检验孔径相容性，不进行成像或源定位。

六层像素化 TES 焦面在像素尺度上提供了这一检验所需的信息。对于第~\ref{subsec:detector_response_zh}~节定义的保留命中，响应模型给出测量能量，事例目录保留 TES 层号和像素位置。位于不同像素或不同层的命中因而同时给出能量分配和有限的空间杆臂，但目录并不直接给出作用先后，也不能分辨像素内部的作用位置。单命中不含康普顿来向信息，因此作为量热谱线候选保留；多命中事例才进入下述孔径相容性检验。

对一个假定的作用次序，设总能量为 $E_0$ 的候选事例在第一个代表位置 $\mathbf r_1$ 沉积能量 $E_1$，随后以能量 $E_{\mathrm{sc}}=E_0-E_1$ 到达第二个位置 $\mathbf r_2$。测量能量给出的康普顿散射角为
\begin{equation}
  \cos\theta_C =
  1 - m_ec^2\left(\frac{1}{E_{\mathrm{sc}}}
  - \frac{1}{E_1+E_{\mathrm{sc}}}\right).
\end{equation}
向量 $\mathbf r_2-\mathbf r_1$ 确定散射后的飞行方向 $\hat{\mathbf u}_{12}$，却不能唯一确定入射方向；它与 $\theta_C$ 共同给出以 $\mathbf r_1$ 为顶点、以 $-\hat{\mathbf u}_{12}$ 为轴的允许入射方向锥面。若该锥面与侧向入口平面内的圆形分析孔径相交，则认为这一作用次序与光学入口相容。质量模型 A 和质量模型 B 的分析孔径半径分别为 $1.898$ 和 $1.900\,\mathrm{cm}$。这些圆盘只定义本层选择所用的几何一致性区域，并不等同于完整物理窗口或准直器开口。

像素内部的位置不确定性在分析中单独处理。目录给出每个命中的分析位置，但不能确定像素内的实际作用点。本文以记录位置为中心，在分析坐标的 $x$、$y$ 和 $z$ 方向分别取 $\pm0.075$、$\pm0.075$ 和 $\pm0.150\,\mathrm{cm}$，构造一个轴对齐代表盒。盒的八个角点和中心给出九个代表位置，而不是九次额外测量。对前两次命中的一个假定次序，两个代表盒形成 $9\times9=81$ 组位置对，也就得到 81 个候选锥。每个锥面沿方位角均匀采样 24 条母线并向入口平面正向延伸；采样交点落入孔径圆盘，或相邻交点之间的线段穿过圆盘，均判为相交。该判据有意从宽：81 个候选锥中只要有一个与圆盘相交，就保留这一作用次序。图~\ref{fig:compton_aperture_consistency} 概括了这一过程。

\begin{figure}[H]
\centering
\includegraphics[width=\linewidth]{figures/fig_compton_aperture_consistency_zh.pdf}
\caption{康普顿锥与分析孔径相容性 veto 示意图。（a）由前两次 TES 命中的测量能量和代表位置构造康普顿锥；锥面给出允许的源方向集合，而不是一条确定的光子轨迹。（b）每个命中使用轴对齐代表盒的八个角点和中心，给定次序形成 $9\times9=81$ 组位置假设。（c）候选锥在入口平面的截线由 24 个方位采样点及其相邻线段近似；任一截线与分析孔径相交即保留，存在有效重建但未相交时否决。无法形成有效锥的事例归入未重建类并保留在率核算中。示意图不按比例。}
\label{fig:compton_aperture_consistency}
\end{figure}

作用次序与像素内部的位置不确定性分别处理。双命中事例检验两种可能的次序，任一次序与孔径相容即予保留。对于至少三命中事例，先用记录的分析位置枚举全部作用次序。设一个候选次序含 $n$ 个命中 $\{(\mathbf r_j,\epsilon_j)\}_{j=1}^{n}$，其康普顿序列残差定义为
\begin{equation}
 Q=\sum_{i=2}^{n-1}\left(\cos\theta_{C,i}-\hat{\mathbf u}_{i-1,i}\!\cdot\!\hat{\mathbf u}_{i,i+1}\right)^2 ,
\end{equation}
其中 $\theta_{C,i}$ 由第 $i$ 个命中的沉积能量 $\epsilon_i$ 及其后的剩余能量 $\sum_{k=i+1}^{n}\epsilon_k$ 计算，$\hat{\mathbf u}_{i-1,i}$ 和 $\hat{\mathbf u}_{i,i+1}$ 分别为相邻两段的飞行方向。无物理解的次序不参与排序。其余次序中取 $Q$ 最小者；若 $Q$ 相同，则选择前两次命中间距较大的次序。随后只对这一最优次序的第一散射锥施加 81 组位置假设检验。

最后的否决规则采取保守口径。只有多命中事例能够按上述规则形成有效重建、且受检候选锥均未与分析孔径相交时，才施加 veto。若所有候选次序均无物理解，则该事例归入未重建类并继续计入最终率，而不是直接拒绝。换言之，不能重建本身不作为该事例属于本底的证据。

\subsection{从第 15 天参考结果到任务时间演化}
\label{sec:mission_time_fold}

探测器输运给出了第 15 天参考环境下的分量响应，但多日浮空期间的残余大气深度、地磁环境和活化库存均随时间变化 \cite{Gallego2025Balloon,Kierans2017}。因此，任务投影不是把第 15 天总本底率作为常数直接延长至 20\,d，而是保留固定质量模型的探测器响应，分别演化聚焦信号、八类瞬发粒子族和延迟活化本底。

本文采用一条 0--20\,d 的合成参考轨迹，共 81 个节点，时间间隔为 0.25\,d。轨迹海拔在 36.25--41.25\,km 之间变化，纬度和经度范围分别为 $33.75^\circ$--$34.25^\circ$\,N 和 $99.75^\circ$--$100.25^\circ$\,E；第 15 天参考状态为 38.75\,km、$34^\circ$\,N 和 $100^\circ$\,E。仪器轴仰角固定为 $45^\circ$。海拔变化同时影响入射粒子场和 511\,keV 大气透过率，经纬度变化用于提供受控的地磁环境调制。该轨迹用于在相同条件下比较质量模型 A 和 B，不代表具体飞行的遥测、源可见性、地球遮挡或重指向历史。

对于给定的大气层顶 511\,keV 线通量，第 15 天以外的聚焦信号率由大气透过率缩放：
\begin{equation}
 S^{\mathrm{dir}}_g(t)
 =S^{\mathrm{dir}}_g(t_{15})
 \frac{T_{\mathrm{atm},511}(t)}{T_{\mathrm{atm},511}(t_{15})}
 =F_{\mathrm{TOA},511}A_{\mathrm{eff},g}T_{\mathrm{atm},511}(t),
 \label{eq:mission_signal_zh}
\end{equation}
其中 $g$ 表示质量模型，$A_{\mathrm{eff},g}$ 已包含探测器响应、主动反符合和 Compton 拓扑/孔径选择的单事例接受度。平面平行近似下，$T_{\mathrm{atm},511}(t)=\exp[-\mu_{\mathrm{air}}X(t)/\sin\epsilon]$。第 15 天节点的垂直剩余大气柱深为 $3.461\,\mathrm{g\,cm^{-2}}$；取 $\epsilon=45^\circ$ 及有效干空气质量衰减系数 $0.08736\,\mathrm{cm^2\,g^{-1}}$ \cite{NISTXCOM}，得到 $T_{\mathrm{atm},511}(t_{15})=0.652034$。任务折叠在每个节点重新计算该透过率。

瞬发本底按入射粒子族分别缩放。对于粒子族 $a$，
\begin{equation}
 R_{\mathrm{prompt},g,a}(t)
 =R_{\mathrm{prompt},g,a}(t_{15})
 \frac{\Phi_a(t)}{\Phi_a(t_{15})}
 =\Phi_a(t)K_{g,a},
 \label{eq:prompt_source_response_zh}
\end{equation}
其中 $\Phi_a(t)$ 为该粒子族在轨迹节点处的环境通量，$K_{g,a}$ 为第 15 天参考输运得到的选后响应系数。七类非 $\gamma$ 瞬发粒子和宽带大气 $\gamma$ 分别使用各自的通量尺度。外推只改变分量归一化；族内能谱、角分布和探测器响应仍采用参考输运模板，因此不需要在 81 个节点重新进行完整粒子输运。

延迟活化不能用第 15 天延迟总率乘以单一比例，因为不同核素具有不同的产生历史和半衰期。由粒子族 $a$ 产生的母核 $k$ 的产生率写为
\begin{equation}
 P_{g,a,k}(t)
 =P_{g,a,k}(t_{15})\frac{\Phi_a(t)}{\Phi_a(t_{15})}
 =\Phi_a(t)Y_{g,a,k},
 \label{eq:activation_production_zh}
\end{equation}
其中 $Y_{g,a,k}$ 由对应质量模型的产生目录确定。相邻节点之间的产生率作线性插值，并与放射性衰变作精确卷积：
\begin{equation}
 \begin{aligned}
 N_{g,a,k}(t_{\ell+1})={}&N_{g,a,k}(t_\ell)e^{-\lambda_k\Delta t_\ell}\\
 &+\int_0^{\Delta t_\ell}P_{g,a,k}(t_\ell+u)
 e^{-\lambda_k(\Delta t_\ell-u)}\,\mathrm du,\\
 A_{g,a,k}(t_\ell)={}&\lambda_kN_{g,a,k}(t_\ell).
 \end{aligned}
 \label{eq:delayed_inventory_zh}
\end{equation}
选后延迟率为
\begin{equation}
 R_{\mathrm{delayed},g}(t)=
 \sum_a\sum_kA_{g,a,k}(t)\varepsilon_{g,a,k}.
 \label{eq:delayed_selected_zh}
\end{equation}
逐“粒子族--母核”推进保留了不同半衰期及选后衰变响应。计算取 $N_{g,a,k}(0)=0$，因此没有计入地面预活化和进入第一个轨迹节点以前的上升段库存。

跨事例分组前的总直接本底率为
\begin{equation}
 B^{\mathrm{dir}}_g(t)
 =\sum_aR_{\mathrm{prompt},g,a}(t)+R_{\mathrm{delayed},g}(t).
 \label{eq:direct_background_time_zh}
\end{equation}
不同物理分量之间还会发生偶然时间符合。本文在第 0、5、10、15 和 20 天设置五个公共时间锚点。在每个锚点，首先根据各分量的瞬时率生成泊松到达时间并合并到同一时间轴，随后执行 $1\,\mu\mathrm{s}$ 事例分组、探测器响应、主动反符合和 Compton 拓扑/孔径选择。锚点计算分别给出本底公共时间修正和聚焦信号存活率：
\begin{equation}
 \widehat\rho_{g,r}=
 \frac{n^{\mathrm{tl}}_{g,r}/T^{\mathrm{rep}}_{g,r}}
 {B^{\mathrm{dir}}_g(t_r)},\qquad
 \widehat\eta_{g,r}=
 \frac{n^{\mathrm{pass}}_{g,r}}{N^{\mathrm{probe}}_{g,r}}.
 \label{eq:anchor_time_factors_zh}
\end{equation}
$\widehat\rho$ 把直接本底率换算为公共时间轴上的最终率；偶然能量合并可能使它略大于 1，因此它不是存活概率。$\widehat\eta$ 才表示弱聚焦信号在相同符合环境中的条件存活率。两者在相邻锚点之间分别作线性插值。质量模型 A、B 的单锚点等效统计曝光分别为 $2\times10^4$ 和 $2\times10^5$\,s。

图~\ref{fig:common_time_anchors} 在施加 $\rho_g(t)$ 以前，用事件级公共时间回放检验未经修正的解析外推。对于每个质量模型，本文在第 0、5、10、15 和 20 天分别重新生成一套泊松到达时间，并执行公共时间分组和全部事例选择，再把回放率与式~\eqref{eq:direct_background_time_zh} 在 81 个节点给出的 $B_g^{\mathrm{dir}}(t)$ 比较。按回放自身的泊松标准误差计，质量模型 A 和 B 的最大偏离分别为 $1.67\sigma$ 和 $1.09\sigma$，没有锚点偏离解析直接率曲线达到 $2\sigma$。因此，在当前轨迹及回放统计精度内，这一结果支持用解析折叠率作为时间外推基线。该比较仍属于闭合检验，而不是五个时刻的独立粒子输运：每次回放使用的分量输入率本身来自式~\eqref{eq:prompt_source_response_zh}--\eqref{eq:delayed_selected_zh}。

\begin{figure}[!ht]
\centering
\includegraphics[width=\linewidth]{figures/section4_current/fig04_time_fold_validation.pdf}
\caption{任务时间本底折叠的事件级闭合检验。（a）实线和虚线为由第 15 天分量响应、逐族环境缩放及逐核素活化演化得到的 81 节点科学窗解析直接本底率；阴影表示有限输运模板的 $1\sigma$ 统计不确定度。由于同一输运模板在所有节点重复使用，该不确定度沿时间相关。空心符号为第 0、5、10、15 和 20 天分别生成的公共时间泊松回放率，误差棒为 $\sqrt{n}/T$；质量模型 A、B 的 $T$ 分别为 $2\times10^4$ 和 $2\times10^5$\,s。（b）回放率与解析率之比，误差棒只包含回放统计误差，水平线表示二者相等。有限输运带与回放误差分别展示，不作合并。该图检验给定分量率下的率映射、公共时间分组和事例选择闭合，不代表五个环境中的独立完整输运。}
\label{fig:common_time_anchors}
\end{figure}

任意时刻的最终信号率和本底率为
\begin{equation}
 S_g(t)=S^{\mathrm{dir}}_g(t)\eta_g(t),\qquad
 B_g(t)=B^{\mathrm{dir}}_g(t)\rho_g(t).
 \label{eq:final_time_rates_zh}
\end{equation}
沿 81 个节点作梯形积分，得到曝光时间 $T_n$ 内的累计计数：
\begin{equation}
 \begin{aligned}
 N_s(T_n)&=\sum_{\ell=1}^{n}\frac{S^{\mathrm{dir}}_g(t_{\ell-1})\eta_g(t_{\ell-1})+S^{\mathrm{dir}}_g(t_\ell)\eta_g(t_\ell)}{2}\,\Delta t_\ell,\\
 N_b(T_n)&=\sum_{\ell=1}^{n}\frac{B^{\mathrm{dir}}_g(t_{\ell-1})\rho_g(t_{\ell-1})+B^{\mathrm{dir}}_g(t_\ell)\rho_g(t_\ell)}{2}\,\Delta t_\ell.
 \end{aligned}
 \label{eq:cumulative_counts_zh}
\end{equation}
式~\eqref{eq:cumulative_counts_zh} 给出的 $N_s(T)$ 和 $N_b(T)$ 是下一节任务灵敏度计算的输入，其累计结果见第~\ref{subsec:mission_fold_validation_zh} 节。

\subsection{任务显著度与最小可分辨线通量}
\label{sec:mission_sensitivity_zh}

对于累计信号计数 $N_s(T)$ 和本底计数 $N_b(T)$，本文采用
\begin{equation}
 Z(T)=\frac{N_s(T)}{\sqrt{N_b(T)}}
 \label{eq:mission_significance_zh}
\end{equation}
定义计数显著度。在本文的高计数、弱信号条件下，该式是泊松计数似然在 $N_s\ll N_b$ 极限下的解析近似 \cite{Cowan2011}。

由于聚焦信号计数与源通量线性相关，而预测本底不随弱源通量变化，曝光时间 $T$ 下的 $q\sigma$ 最小可分辨大气层顶线通量可直接写为
\begin{equation}
 F_{q\sigma}(T)=\fref\frac{q}{Z(T;\fref)},
 \qquad q\in\{3,5\}.
 \label{eq:minimum_flux_zh}
\end{equation}
达到给定显著度所需的观测时间由累计曲线 $Z(T)=q$ 的首次交点确定。

81 个任务节点重复使用同一批输运事例，因此各节点的统计误差并不独立。本文先累计每条事例在全部节点上的带权贡献，再计算任务累计计数的统计误差。五个时间锚点、有效面积和信号探针的统计误差也一并传播到时间曲线和流量阈值；阈值到达时间在相邻节点之间线性插值。


\section{结果}
\label{sec:results}

本节沿图~\ref{fig:workflow} 的端到端流程组织结果：首先建立质量模型 A 的本底与响应基线，并由最终入选事例追溯限制本底的材料、母核素和空间区域；随后检验质量模型 B 的完整侧烟囱构型，最后在共同时间轴上比较两套设计的任务性能。两模型采用相同的能量响应、科学窗、到达时间分组和 Compton 拓扑/孔径判据；主动屏蔽的通道构成属于各自完整设计的一部分。\revblue{为诊断最终本底组成，本节把同一八族瞬发分支进一步表示为宽带大气 $\gamma$ 分量和七类非 $\gamma$ 粒子的合并分量；这一报告拆分不改变源构建或率归一。}

\subsection{质量模型 A 的基线本底与选择性能}
\label{subsec:model_a_baseline_zh}
\label{subsec:reference_background_results_zh}

质量模型 A 构成本文的基线。表~\ref{tab:phase2_cutflow} 给出第 15 天直接、无符合样本在 $[510.58,511.42)\keV$ 科学窗内的完整 cut-flow。主动 BGO 门将总本底由 3086 条记录、$5.43630\cps$ 降至 443 条记录、$6.07099\times10^{-2}\cps$，相当于去除预 veto 率的 $98.883\%$。\revblue{七类非 $\gamma$ 瞬发粒子}由 $2.91217\cps$ 降为零；宽带大气 $\gamma$ 本底由 $2.35667\cps$ 降至 $1.78923\times10^{-2}\cps$；延迟活化由 $0.167460\cps$ 降至 $4.28176\times10^{-2}\cps$。

最终 Compton 拓扑/孔径选择后，本底降至 400 条记录、$5.30571\times10^{-2}\cps$，即保留主动门后带权率的 $87.394\%$。同一 $37{,}194$ 光子焦面样本中，测量窗内的 28,612 条孤立聚焦信号记录全部通过主动门；最终保留 28,006 条，对应有效面积 $15.12324\pm0.04492\,\mathrm{cm^2}$，为测量窗信号的 $97.882\%$。主动反符合承担主要的率压低，拓扑检验移除较小的残余，同时保留绝大多数聚焦信号。

\begin{table}[H]
\centering
\caption{质量模型 A 科学窗 cut-flow。本底为第 15 天直接、无符合记录数/率；信号为共同 $37{,}194$ 光子焦面样本的记录数/有效面积。}
\label{tab:phase2_cutflow}
\footnotesize
\begin{tabular}{lccc}
\toprule
计算流 & 预 veto & 主动 veto 后 & 最终拓扑/孔径后 \\
\midrule
\revblue{七类非 $\gamma$ 瞬发粒子} & 869 / $2.91217\cps$ & 0 / 0 & 0 / 0 \\
宽带大气 $\gamma$ 本底 & 922 / $2.35667\cps$ & 7 / $1.78923\times10^{-2}\cps$ & 6 / $1.53363\times10^{-2}\cps$ \\
延迟活化 & 1295 / $0.167460\cps$ & 436 / $4.28176\times10^{-2}\cps$ & 394 / $3.77209\times10^{-2}\cps$ \\
总本底 & 3086 / $5.43630\cps$ & 443 / $6.07099\times10^{-2}\cps$ & 400 / $5.30571\times10^{-2}\cps$ \\
聚焦信号 & 28612 / $15.45048\,\mathrm{cm^2}$ & 28612 / $15.45048\,\mathrm{cm^2}$ & 28006 / $15.12324\,\mathrm{cm^2}$ \\
\bottomrule
\end{tabular}
\end{table}

最终直接本底的非零中心值由延迟活化和宽带大气 $\gamma$ 本底构成（表~\ref{tab:reference_background_budget_zh}）。延迟活化为 $3.77209\times10^{-2}\cps$，占 $71.095\%$；宽带大气 $\gamma$ 本底为 $1.53363\times10^{-2}\cps$，占 $28.905\%$；\revblue{七类非 $\gamma$ 瞬发粒子}的中心值为零，但有限样本下不能将该中心值解释为不存在瞬发本底。各分量的单事例权重不同，因此记录数不能代替有效样本量；总本底 400 条记录对应 $N_{\mathrm{eff}}=46.77$。

\begin{table}[H]
\centering
\caption{质量模型 A 第 15 天科学窗最终直接本底。标准误差和有效样本量由不等权事例的 $\sum_i w_i^2$ 计算。}
\label{tab:reference_background_budget_zh}
\footnotesize
\begin{tabular}{lrrrrr}
\toprule
分量 & 记录数 & 率 / $\cps$ & 标准误差 / $\cps$ & $N_{\mathrm{eff}}$ & 总本底 / \% \\
\midrule
延迟活化 & 394 & $3.77209\times10^{-2}$ & $4.58152\times10^{-3}$ & 67.79 & 71.095 \\
宽带大气 $\gamma$ 本底 & 6 & $1.53363\times10^{-2}$ & $6.26100\times10^{-3}$ & 6.00 & 28.905 \\
\midrule
总计 & 400 & $5.30571\times10^{-2}$ & $7.75825\times10^{-3}$ & 46.77 & 100.000 \\
\bottomrule
\end{tabular}
\end{table}

\subsection{质量模型 A 的本底来源与空间耦合}
\label{subsec:mass_model_a_origins_zh}

总率回答本底量级，但不能说明事例为何通过最终选择。质量模型 A 的入选目录因此按初级入射粒子族、放射性母核的产生材料、源体积--母核素组合以及相对 TES 的空间区域逐层追溯。

按初级入射粒子族分解，质子诱发活化是最大的单项来源，最终率为 $2.48463\times10^{-2}\cps$（35 条记录），占质量模型 A 总直接本底的 $46.829\%$。其余延迟分量依次为 $\alpha$ 诱发的 $7.49239\times10^{-3}\cps$（26 条，$14.121\%$）、中子诱发的 $3.60054\times10^{-3}\cps$（11 条，$6.786\%$）、$\gamma$ 诱发的 $1.68617\times10^{-3}\cps$（260 条，$3.178\%$）、$e^+$ 诱发的 $7.02083\times10^{-5}\cps$（24 条，$0.132\%$）和 $e^-$ 诱发的 $2.52755\times10^{-5}\cps$（38 条，$0.048\%$）。这些分量之和为质量模型 A 的最终延迟率 $3.77209\times10^{-2}\cps$。

按材料分类，铜贡献带权选后延迟率的 $70.87\%$。具体到源体积--母核素组合，主要具名来源是 L2 铜支撑环中的 $^{62}$Cu、50\,mK 混合室铜板中的 $^{62}$Cu 和 L0 TES 热沉环中的 $^{62}$Cu；其后的具名组包括 Bi 上半圆柱中的 $^{199}$Po，以及 50\,mK 混合室铜板和 L2 铜支撑环中的 $^{64}$Cu。这里的排序已经包含输运、响应和全部选后条件，不能由材料总质量或总活度直接推出。

空间分解说明了几何耦合：DR 混合室及多级冷盘贡献带权选后延迟率的 $32.18\%$，TES 近场结构贡献 $47.29\%$，两者合计 $79.47\%$。按最终延迟率换算，前者约为 $1.21398\times10^{-2}\cps$，后者约为 $1.78397\times10^{-2}\cps$。质量模型 A 的主要延迟负担并非铜总量本身，而是部分活化铜及近场结构仍与 TES 入选条件保持较强的空间耦合。

这一来源诊断提出了一个明确但尚非因果证明的设计问题：若改变焦面与冷端结构的相对位置，并用局部主动屏蔽重新包围探测器，能否在保留聚焦信号响应的同时削弱这部分选后耦合？质量模型 B 是对该完整设计问题的独立实现。由于它同时改变物理入口、局部 W 结构、BGO 布局和部分材料清单，后文只能评价这套完整构型的条件性能，不能把任一差值单独归因于侧烟囱或焦面位移。

\FloatBarrier

\subsection{由来源诊断到质量模型 B 的完整构型}
\label{subsec:optimized_shield_results_zh}

质量模型 B 将六层 TES 沿局部 $z'=-2.80$\,cm 的侧烟囱轴布置。烟囱 Al 壁的内、外半径分别为 10.75 和 11.05\,cm，壁厚为 3\,mm，并与稀释制冷机外壳形成连续接口；烟囱底部与稀释制冷机底部之间保留 0.25\,cm 间隙。焦面由外、内方孔分别为 6.00 和 5.40\,cm、轴向深度为 2.00\,cm 的四边 W 框支撑。三个主动 BGO 体积覆盖烟囱侧面、前端光学环和后端冷孔环。主动门仍按同一 $1\,\mu\mathrm{s}$ 公共时间组累计全部 BGO 沉积，并要求组内 BGO 总沉积严格低于 50\,keV。

\revmark{18}图~\ref{fig:optimized_shield_background_zh} 以相同仪器坐标和比例比较两套冷端结构，并在局部放大图中显示 TES、主要铜支撑、W 孔径/框架和主动 BGO 的相对位置；表~\ref{tab:final_geometry_zh} 列出输运采用的质量模型 B 尺寸。

\revdeleted{R18：A/B 有效面积比较移至质量模型 B 的全链结果，不在构型介绍中提前报告。}

\begin{figure}[!ht]
\centering
\includegraphics[width=\linewidth]{figures/section4_current/fig09_geometry_response_detailed.pdf}
\caption{质量模型 A 与质量模型 B 的冷端--焦面 $x'$--$z'$ 切面及选后延迟源位置。（a）、（b）采用相同坐标范围和比例，显示多级铜冷盘、六层 TES、主动 BGO、W 孔径/框架以及质量模型 B 的侧烟囱和铜冷指；（c）、（d）以相同物理尺度放大两个焦面邻域，并在质量模型 A 中标出 L2 铜支撑环、L0 铜热沉环和 50\,mK MXC 铜盘。每个符号对应一条第 15 天最终入选延迟事例的母核产生位置，颜色表示源材料，点面积按带权率采用两模型共同标度并设置最小可见尺寸。切面按输运几何尺寸绘制；为清晰起见，未显示小型 CSG 开孔和服务件。}
\label{fig:optimized_shield_background_zh}
\end{figure}

\begin{table}[!ht]
\centering
\caption{质量模型 B 侧烟囱在局部仪器坐标中的几何参数。}
\label{tab:final_geometry_zh}
\footnotesize
\begin{tabular}{p{0.23\linewidth}p{0.34\linewidth}p{0.34\linewidth}}
\toprule
部件 & 最终设置 & 设计作用 \\
\midrule
Al 烟囱 & 内/外半径 10.75/11.05 cm；壁厚 3 mm & 将 TES 横向布置于多级冷端区域之外 \\
共同轴与间隙 & $z'=-2.80$ cm；烟囱/DR 底部 $-13.85/-14.10$ cm & 连续接口并保留 0.25 cm 间隙 \\
TES 焦面 & 沿烟囱轴布置六层 Ta 像素 & 保持探测器响应和多层阻止深度 \\
W 焦面框 & 外/内方孔 6.00/5.40 cm；深 2.00 cm & 紧凑局部支撑及孔径定义 \\
冷端/光学孔 & 1.50 cm 冷孔；6.00 cm 上端光学开孔 & 保留低温服务和聚焦光束通道 \\
主动 BGO & 侧屏、前端光学环、后端冷孔环 & 环绕烟囱并按共时间组累计的反符合 \\
BGO 阈值 & 50 keV 离线沉积能 & 共同主动 veto 判据 \\
\bottomrule
\end{tabular}
\end{table}

\revblue{后文报告的有效面积比较只约束聚焦光子样本}，并不意味着两套几何具有相同的弥散光子接受度。质量模型 A 使用半径约 1.898\,cm 的入口包络和 624 孔 W 准直结构；质量模型 B 使用半径 2.70\,cm 的物理窗口和开放四边 W 框，主动屏蔽构型也不同。因此，本节比较的是共同源、响应、计时和选择约定下的两套完整设计。质量模型 B 的本底和灵敏度以其入口、W 框、BGO 与主动通道组合为条件；现有数据没有提供仅改变侧烟囱、同时冻结其余结构的消融样本。

\FloatBarrier

\subsection{质量模型 B 的全链本底与来源重分配}
\label{subsec:optimized_fullchain_zh}

质量模型 B 的探测器侧事例目录沿用上述 420\,eV FWHM 的 TES 响应、$[510.58,511.42)\keV$ 科学窗、主动 BGO 门以及同一 Compton 拓扑/孔径检验。表~\ref{tab:optimized_cutflow_zh} 给出第 15 天直接、无符合样本的完整 cut-flow。

主动 BGO 将总本底由 7684 条记录、$2.38372\cps$ 降至 137 条记录、$1.11379\times10^{-2}\cps$，即只剩预 veto 率的 $0.4672\%$。测量窗内的 28,434 条孤立聚焦信号记录全部通过主动门。随后，Compton 拓扑/孔径检验将本底由 137 条降至 126 条，保留主动门后带权率的 $96.492\%$；聚焦信号由 28,434 条降至 27,936 条，对应最终有效面积 $15.08544\pm0.04503\,\mathrm{cm^2}$。主动反符合承担主要的率压低，最终拓扑检验移除较小的残余，同时保留 $98.249\%$ 的测量窗聚焦信号。\revblue{在共同的轴上信号样本和选择框架下，质量模型 B 的最终有效面积为质量模型 A 的 $99.75\%$。}

\begin{table}[H]
\centering
\caption{质量模型 B 经响应卷积的科学窗 cut-flow。本底为第 15 天直接、无符合记录数/率；信号为记录数/有效面积。}
\label{tab:optimized_cutflow_zh}
\footnotesize
\begin{tabular}{lccc}
\toprule
计算流 & 预 veto & 主动 BGO 后 & 最终拓扑/孔径后 \\
\midrule
\revblue{七类非 $\gamma$ 瞬发粒子} & 1448 / $0.878692\cps$ & 3 / $1.81131\times10^{-3}\cps$ & 3 / $1.81131\times10^{-3}\cps$ \\
宽带大气 $\gamma$ 本底 & 2987 / $1.37944\cps$ & 12 / $5.35535\times10^{-3}\cps$ & 12 / $5.35535\times10^{-3}\cps$ \\
延迟活化 & 3249 / $0.125586\cps$ & 122 / $3.97120\times10^{-3}\cps$ & 111 / $3.58054\times10^{-3}\cps$ \\
总本底 & 7684 / $2.38372\cps$ & 137 / $1.11379\times10^{-2}\cps$ & 126 / $1.07472\times10^{-2}\cps$ \\
聚焦信号 & 28434 / $15.35436\,\mathrm{cm^2}$ & 28434 / $15.35436\,\mathrm{cm^2}$ & 27936 / $15.08544\,\mathrm{cm^2}$ \\
\bottomrule
\end{tabular}
\end{table}

质量模型 B 的最终直接本底由\revblue{七类非 $\gamma$ 瞬发粒子} $1.81131\times10^{-3}\cps$、宽带大气 $\gamma$ 本底 $5.35535\times10^{-3}\cps$ 和延迟活化 $3.58054\times10^{-3}\cps$ 构成，分别占总率的 $16.854\%$、$49.830\%$ 和 $33.316\%$；总直接率为 $1.07472\times10^{-2}\cps$。与质量模型 A 相比，延迟活化和宽带大气 $\gamma$ 分量分别降至其 $9.492\%$ 和 $34.920\%$，总率降至 $20.256\%$，即降低 4.94 倍。质量模型 A 的\revblue{七类非 $\gamma$ 瞬发粒子}中心值为零，因而不能用中心值为这项构造跨模型比值。

三项本底的相对权重随完整构型改变：质量模型 A 由延迟活化主导，而质量模型 B 中宽带大气 $\gamma$ 本底成为最大分量，并出现由三个正电子族模板支撑的瞬发中心值。这是有限统计下的来源重分配，而不是对正电子本底的稳健识别。由于入口、W 框、BGO 与局部材料同时变化，A/B 对照不能把总率或组成差异唯一归因于其中某个结构。

\begin{table}[!ht]
\centering
\caption{质量模型 B 第 15 天最终直接本底的有限输运支撑。标准误差和有效样本量由不等权事例的 $\sum_i w_i^2$ 计算。}
\label{tab:optimized_component_precision_zh}
\scriptsize
\resizebox{\linewidth}{!}{%
\begin{tabular}{lrr@{\hspace{1.4em}}rrr}
\toprule
分量 & 记录数 & 率 / $\cps$ & 标准误差 / $\cps$ & $N_{\mathrm{eff}}$ & 总本底 / \% \\
\midrule
\revblue{七类非 $\gamma$ 瞬发粒子} & 3 & $1.81131\times10^{-3}$ & $1.04576\times10^{-3}$ & 3.00 & 16.854 \\
宽带大气 $\gamma$ 本底 & 12 & $5.35535\times10^{-3}$ & $1.55631\times10^{-3}$ & 11.84 & 49.830 \\
延迟活化 & 111 & $3.58054\times10^{-3}$ & $5.36971\times10^{-4}$ & 44.46 & 33.316 \\
\midrule
总本底 & 126 & $1.07472\times10^{-2}$ & $1.95040\times10^{-3}$ & 30.36 & 100.000 \\
\bottomrule
\end{tabular}}
\end{table}

表~\ref{tab:optimized_component_precision_zh} 给出三项分量的不等权统计支撑。最终样本共 126 条记录，对应 $N_{\mathrm{eff}}=30.36$。\revblue{七类非 $\gamma$ 瞬发粒子}的三个入选模板均来自正电子族，但其有限样本区间较宽；这一中心值只用于总率核算，不足以确认稳定的正电子主导分量。质量模型 A 的\revblue{七类非 $\gamma$ 瞬发粒子}中心值为零，同样不能解释为该类本底不存在。

质量模型 B 的铜来源为 $(2.26708\pm0.42826)\times10^{-3}\cps$，占其最终延迟率的 $63.32\%$。DR/MXC 及多级冷盘来源由质量模型 A 的 $(1.21398\pm0.26869)\times10^{-2}\cps$ 降至质量模型 B 的 $(1.40604\pm1.09289)\times10^{-4}\cps$；质量模型 B 的该分量仅有 7 条来源记录，因而相对统计误差较大。总延迟率则由 $(3.77209\pm0.45815)\times10^{-2}$ 降至 $(3.58054\pm0.53697)\times10^{-3}\cps$。质量模型 B 的主要具名组仍包括 L5 铜环和 TES 热沉中的 $^{62}$Cu。图~\ref{fig:geometry_response_rates_zh} 汇总了这种空间重分配及完整构型的响应比；选后铜贡献因而是被削弱并重新分配，而不是在材料层面被完全消除。

\begin{figure}[!ht]
\centering
\includegraphics[width=\linewidth]{figures/section4_current/fig09b_geometry_response_rates.pdf}
\caption{选后延迟本底的空间重分配与两套完整设计的响应比。（a）质量模型 A、B 的第 15 天最终延迟率按冷端、TES 近场和其他结构分解。（b）质量模型 B/A 的冷端延迟、铜来源延迟、总延迟、科学窗总本底和聚焦信号有效面积比；虚线表示两模型相等。误差棒由不等权事例的二阶矩给出，跨模型比值按两项独立传播；不包含共同时间或物理系统误差。由于两模型的入口、局部 W 结构、BGO 布局、近场材料和主动通道构成不同，该图比较完整设计的条件响应，而不是单因素消融。}
\label{fig:geometry_response_rates_zh}
\end{figure}

图~\ref{fig:day15_selection_spectra_zh} 给出的 480--550\,keV 宽窗谱提供独立于窄窗的选择诊断。质量模型 A 的宽窗带权率由预 veto 的 $10.7984\cps$ 降至联合主动门后的 $0.140797\cps$，最终为 $0.121285\cps$；质量模型 B 相应由 $4.92157\cps$ 降至 $0.0735321\cps$，最终为 $0.0655197\cps$。这与科学窗 cut-flow 给出相同的阶段排序：主动反符合承担主要抑制，Compton 拓扑/孔径检验对残余作较小修正。

\begin{figure}[!ht]
\centering
\includegraphics[width=\linewidth]{figures/section4_current/fig05_day15_selection_spectra_current.pdf}
\caption{第 15 天质量模型 A、B 的 TES 测量能量微分计数率谱。（a）、（b）为 480--550\,keV 宽带，0.25\,keV 分析能箱按四箱合并为 1\,keV；（c）、（d）保留 508--514\,keV 内原生 0.25\,keV 分箱，灰带及虚线标出 $[510.58,511.42)\keV$ 科学窗。三条曲线依次为主动反符合前、主动反符合后和最终 Compton 拓扑/孔径选择后；宽带面板的半透明带和局部面板的逐箱误差棒表示由 $\sqrt{\sum_i w_i^2}$ 得到的有限输运 $1\sigma$ 误差。零权重分箱保持断开，不作平滑或伪计数处理。纵轴为输运和 TES 响应后的探测器微分计数率。}
\label{fig:day15_selection_spectra_zh}
\end{figure}

质量模型 B 的延迟活化与宽带大气 $\gamma$ 本底中心率均低于质量模型 A，最终组成由质量模型 A 的延迟活化主导转为质量模型 B 的宽带大气 $\gamma$ 本底占比最高。该比较支持对质量模型 B 完整构型作条件性能评价，但不构成侧烟囱这一单一几何因素的因果测量。

\FloatBarrier

\subsection{共同时间轴折叠与累计计数}
\label{subsec:mission_fold_validation_zh}

在第 0、5、10、15 和 20 天的五个锚点，各物理分量重新生成泊松到达并在共同时间轴上完成分组、响应、主动反符合和 Compton 选择。共同时间最终率与直接率之比定义本底修正 $\widehat\rho_g$，聚焦信号重放给出偶然存活率 $\widehat\eta_g$；二者线性插值到 81 个任务节点。质量模型 A 的 $\widehat\rho_g$ 为 0.9307--1.0060，质量模型 B 为 0.9901--1.0293（图~\ref{fig:common_time_anchors}）。

如图~\ref{fig:common_time_anchors} 所示，五个事件级回放率在回放自身的 $2\sigma$ 统计范围内均与相应解析直接率曲线相容。这一比较建立了公共时间实现的数值闭合并量化修正 $\widehat\rho_g$，但不能替代非参考环境中的独立粒子输运。

聚焦信号的 $\widehat\eta_g$ 在质量模型 A 中为 0.96810--0.97039，在质量模型 B 中为 0.99547--0.99580；第 15 天分别为 0.9682365 和 0.995472。质量模型 B 的较高值来自其较低的符合占有率，而非不同的信号选择。$\rho_g$ 和 $\eta_g$ 插值到 0--20\,d、步长 0.25\,d 的 81 个节点后，与同一轨迹透过率和分量响应作梯形积分。

\revblue{将第 15 天的条件存活率与各模型的孤立事例最终有效面积及该节点的大气透过率结合，在第~\ref{sec:requirements} 节定义的参考通量下，质量模型 A 和质量模型 B 的最终信号率分别为 $2.29144\times10^{-3}$ 和 $2.35000\times10^{-3}\cps$。质量模型 B 的孤立事例有效面积略低，但偶然符合损失较小，因而其最终参考信号率略高。}

20\,d 时，质量模型 A 累计 3928.83 个参考流量信号和 87632.34 个本底计数；质量模型 B 累计 4028.15 个信号和 18568.89 个本底计数。宽带大气 $\gamma$ 与其余本底的累计计数在质量模型 A 中分别为 25760.39 和 61871.95，在质量模型 B 中分别为 9359.44 和 9209.45；后者同时包含\revblue{七类非 $\gamma$ 瞬发粒子}和延迟活化，不能整体解释为活化。质量模型 B 的累计本底为质量模型 A 的 $21.19\%$。

\subsection{任务显著度与最小可分辨线通量}
\label{subsec:final_mission_performance_zh}

图~\ref{fig:optimized_mission_significance_zh} 和表~\ref{tab:primary_sensitivity_zh} 汇总两套设计的任务性能。在 $\fref$ 下，20\,d 计数显著度在质量模型 A、B 中分别为 13.272 和 29.561。相应的大气层顶 $3\sigma$ 最小可分辨通量分别为 $(5.425\pm0.403)\times10^{-5}$ 和 $(2.436\pm0.223)\times10^{-5}\phcms$。这些误差为局部传播 $1\sigma$，不包含完整联合不确定度。

\begin{figure}[!ht]
\centering
\includegraphics[width=0.98\linewidth]{figures/section4_reference_flux_20260830/fig10_mission_performance_gaussian.pdf}
\caption{质量模型 A 与质量模型 B 在共同分析约定下的 81 节点条件任务性能。（a）大气层顶参考通量 $\fref=2.4\times10^{-4}\phcms$ 下的计数显著度，水平线标出 $3\sigma$ 与 $5\sigma$；（b）相应的 $3\sigma$ 大气层顶最小可分辨通量。两条曲线采用相同源、探测器响应、科学窗、计时、Compton 拓扑/孔径及大气透过约定；入口、局部 W 结构和主动通道按各自完整设计定义。}
\label{fig:optimized_mission_significance_zh}
\end{figure}

\begin{table}[!ht]
\centering
\caption{质量模型 A 与质量模型 B 在共同分析合同下的条件任务投影；流量阈值后的误差为局部传播的 $1\sigma$ 统计不确定度，不是联合置信区间。}
\label{tab:primary_sensitivity_zh}
\footnotesize
\resizebox{\linewidth}{!}{%
\begin{tabular}{lrr}
\toprule
量 & 质量模型 A & 质量模型 B \\
\midrule
选后有效面积 / $\mathrm{cm^2}$ & $15.12324\pm0.04492$ & $15.08544\pm0.04503$ \\
第 15 天直接本底 / $\cps$ & $5.30571\times10^{-2}$ & $1.07472\times10^{-2}$ \\
第 15 天 $\fref$ 信号 / $\cps$ & $2.29144\times10^{-3}$ & $2.35000\times10^{-3}$ \\
20 d 本底计数 & 87632.34 & 18568.89 \\
20 d $\fref$ 信号计数 & 3928.83 & 4028.15 \\
$Z_{20\mathrm d}=N_s/\sqrt{N_b}$ & 13.272 & 29.561 \\
$\fref$ 下的 $T_{3\sigma}$ / d & 0.669 & 0.198 \\
$\fref$ 下的 $T_{5\sigma}$ / d & 1.913 & 0.452 \\
$F_{3\sigma}(20\,\mathrm d)$ / $\phcms$ &
$(5.425\pm0.403)\times10^{-5}$ &
$(2.436\pm0.223)\times10^{-5}$ \\
$F_{5\sigma}(20\,\mathrm d)$ / $\phcms$ &
$(9.042\pm0.671)\times10^{-5}$ &
$(4.059\pm0.372)\times10^{-5}$ \\
\bottomrule
\end{tabular}}
\end{table}

\FloatBarrier

流量阈值的变化可以直接拆成累计本底和累计信号核两项。质量模型 B 的 20\,d 本底为质量模型 A 的 $21.19\%$，而其 20\,d 累计信号核为质量模型 A 的 1.025278 倍，因此
\begin{equation}
 \frac{F_{\min,A}}{F_{\min,B}}
 =\sqrt{\frac{N_{b,A}}{N_{b,B}}}\frac{K_B^{20\mathrm d}}{K_A^{20\mathrm d}}
 =2.172397\times1.025278
 =2.227311 .
 \label{eq:fmin_ratio_decomposition_zh}
\end{equation}
这里的 $K_g^{20\mathrm d}$ 是 81 节点积分后的累计信号核。式~\eqref{eq:fmin_ratio_decomposition_zh} 表明，2.227 倍的流量阈值差异主要来自总本底的平方根项，并由累计信号核差异作进一步修正。

质量模型 B 保留质量模型 A 的 $99.75\%$ 聚焦有效面积，并将 20\,d 累计本底和 $3\sigma$ 流量阈值分别降至 $21.19\%$ 和 $44.90\%$。由于入口、W 框、主动屏蔽及局部材料同时变化，这一结果是两套完整设计的条件比较，不能识别单一几何因素。

\FloatBarrier

\section{轨道与月面环境的条件响应转移}
\label{sec:environment_screening}

为判断质量模型 B 的残余本底对辐射环境变化的响应，本文构造了一个探索性转移计算。目标源场包括 530\,km 近赤道、南大西洋异常区（SAA）外的 LEO 代理，暴露月面代理，以及两个太阳调制端点的安静日--地 L2 环境。LEO 入射谱采用低地球轨道 $\gamma$ 射线本底计算中的分量定义 \cite{Cumani2019LEO}。月面源场将 REDMoon 在 2019 年太阳极小期给出的 Apollo-17 月壤上行谱与月面可见宇宙 $\gamma$ 天空拼接 \cite{Dobynde2021REDMoon}。日--地 L2 源集包含 $\gamma$、质子、$\alpha$、电子和正电子；其中质子力场调制势在 2009 年太阳极小/GCR 高通量端点取 $0.3793$\,GV，在 2014 年太阳极大/GCR 低通量端点取 $0.803$\,GV，$\alpha$ 分量采用相应的长期 GCR 参数化 \cite{Lotti2021Athena,Usoskin2005,Kuznetsov2017}。能量均表示单个入射粒子的总动能，其中 $\alpha$ 能量按整颗原子核而非每核子定义。

保留的质量模型 B 转移响应把每条瞬发候选追溯至入射粒子族和初级能量；延迟候选则与具有相同入射族、母核状态和产生体积的带权活化产生能量分布耦合。记 $e=0$ 为 38\,km 气球参考环境，则响应流 $x\in\{\gamma,\mathrm d\}$ 在环境 $e$ 中的表示率为
\begin{equation}
 \widetilde B_{x,e}=
 \sum_{i\in\mathcal R_x}w_{i,0}
 \int p_i(E)
 \frac{\phi_{a_i,e}(E)}{\phi_{a_i,0}(E)}
 \,\mathrm dE ,
 \label{eq:environment_reweight}
\end{equation}
其中 $w_{i,0}$ 为气球参考的选后权重，$a_i$ 为入射粒子族，$\phi_{a,e}$ 为按该源声明角域积分的能谱，$p_i(E)$ 为归一化初级能量核。对瞬发 $\gamma$ 记录，$p_i(E)$ 集中在该记录的输运初级能量；对延迟记录，它表示相应的活化产生能量分布。只有目标谱与参考谱共同具有非零支撑的能区才进入转移。气球响应的方向模板保持不变，因此式~\eqref{eq:environment_reweight} 是重要性重加权筛查，而不是新的方向匹配输运。

为使环境转移与第~\ref{sec:results}~节的质量模型 B 本底预算采用同一归一化口径，两个可表示流分别锚定到第 15 天中心率：
\[
 B_{\gamma,0}=5.355349\times10^{-3}\cps,\qquad
 B_{\mathrm d,0}=3.580538\times10^{-3}\cps .
\]
定义
\[
 c_\gamma=\frac{B_{\gamma,0}}{\widetilde B_{\gamma,0}}=0.924290,
 \qquad
 c_{\mathrm d}=\frac{B_{\mathrm d,0}}{\widetilde B_{\mathrm d,0}}=0.987825,
\]
则两流表示本底及其环境比为
\begin{equation}
 B^{(2)}_e=c_\gamma\widetilde B_{\gamma,e}
          +c_{\mathrm d}\widetilde B_{\mathrm d,e},
 \qquad
 {\cal R}^{(2)}_e=\frac{B^{(2)}_e}{B^{(2)}_0}.
 \label{eq:environment_two_stream}
\end{equation}
这是流总量层面的重新锚定近似：它使当前两个流的总率闭合，但不等价于证明保留转移样本内每条事例与当前事例具有相同的流内相对权重。

当前\revblue{七类非 $\gamma$ 瞬发粒子的合并残余}为 $1.811307\times10^{-3}\cps$，占完整气球率的 16.854\%，由 3 条正电子记录支撑。紧凑转移响应没有瞬发正电子的初级能量核，因此无法为这一项推导族分辨的目标环境率。为构造以正式全本底气球阈值为锚点的无量纲指数，本文作出一个明确的闭合约定：完整本底的相对变化由 $ {\cal R}^{(2)}_e$ 表示；等价地，未映射的瞬发份额只在式~\eqref{eq:environment_flux_screen} 中随两流总比变化，但不作为独立转移得到的正电子率报告。该约定是筛查指数最主要的未映射系统误差，所得数值并非目标环境完整本底预测。

\revblue{对相同的 20\,d 曝光和最终有效面积，记 $T_{20\mathrm d}=20\,\mathrm d=1.728\times10^6\,\mathrm{s}$，正式气球信号核与无衰减、连续在源信号核之比为}
\begin{equation}
 \tau_{\mathrm B}=
 \frac{K_{\mathrm B}^{20\mathrm d}}
 {A_{\mathrm{eff,B}}\revblue{T_{20\mathrm d}}}=0.643861 .
 \label{eq:environment_signal_factor}
\end{equation}
非气球入射平面取 $\tau_e=1$，并以正式气球阈值 $F_{3\sigma,\mathrm B}=2.435681\times10^{-5}\phcms$ 为锚点，定义
\begin{equation}
 \frac{F_{3\sigma,e}^{\mathrm{screen}}}{F_{3\sigma,\mathrm B}}=
 \begin{cases}
  1, & e=0,\\[3pt]
  (\tau_{\mathrm B}/\tau_e)\sqrt{{\cal R}^{(2)}_e}, & e\ne0 .
 \end{cases}
 \label{eq:environment_flux_screen}
\end{equation}
气球锚点是大气层顶流量，非气球数值则对应当地仪器入射平面。权重集中度 $N_{\mathrm{eff},w}=(\sum_iw_i)^2/\sum_iw_i^2$ 与正文正式总本底有效样本量分开报告。

质量模型 B 的选后事例把窄科学窗与更宽的初级能量区间联系起来。响应加权的 10--90\% 区间分别为：瞬发宽带 $\gamma$ 的 $0.598$--$29.6$\,MeV，以及延迟 $\gamma$、正电子、质子、$\alpha$ 和中子产生的 $0.0745$--$611$\,GeV、$0.115$--$2.01$\,GeV、$1.37$--$48.4$\,GeV、$18.7$--$82.9$\,GeV 和 $0.0197$--$195$\,GeV。$\alpha$ 区间按每个原子核总动能给出。这些响应能区构成式~\eqref{eq:environment_reweight} 中环境能谱加权的主要支撑。

\begin{figure}[!ht]
\centering
\includegraphics[width=\linewidth]{figures/environment_screening_20260830/fig11_environment_transfer_zh.pdf}
\caption{质量模型 B 的响应加权环境筛查。（a）按各源声明角域积分的微分 $\gamma$ 源谱；阴影表示瞬发 $\gamma$ 选后响应的初级能量 10--90\% 区间。（b）质子谱，包括日--地 L2 的 2009 年 GCR 高通量和 2014 年 GCR 低通量端点；阴影为相应延迟响应区间。阴影表示响应支撑而非源不确定度。各源保留自身物理角域，曲线不是共同各向同性强度的测量。（c）由式~\eqref{eq:environment_flux_screen} 得到的两流 $3\sigma$ 筛查指数；数值标注单位为 $10^{-5}\phcms$。空心点强调非气球数值是响应转移指数，不是匹配输运灵敏度或置信区间。}
\label{fig:environment_transfer}
\end{figure}

表~\ref{tab:environment_screening} 给出中心筛查指数。安静非 SAA LEO 源场的 $ {\cal R}^{(2)}=0.5627$；去除气球大气对信号核的衰减后，阈值指数为 0.4830，即 $1.176\times10^{-5}\phcms$。这一表观优势只适用于所设定的静态非 SAA 代理，不包含俘获粒子，也不包含穿越 SAA 时产生并在离开后继续衰变的活化库存。

暴露月面代理的 $ {\cal R}^{(2)}$ 则为 23.02，主要由瞬发响应能区内的月壤相关 $\gamma$ 场驱动。无大气衰减部分抵消了本底增加，所得筛查指数为 3.089，对应 $F_{3\sigma}^{\mathrm{screen}}=7.525\times10^{-5}\phcms$。该数值属于所采用的暴露月面源场和孤立探测器--低温系统，并不代表着陆器或月面基地构型。

\begin{table}[!ht]
\centering
\caption{质量模型 B 的两流环境筛查。表示本底仅包含宽带 $\gamma$ 与延迟活化。$F_{3\sigma}^{\mathrm{screen}}$ 通过式~\eqref{eq:environment_flux_screen} 归一到当前正式气球阈值；气球值位于大气层顶，其余值位于当地入射平面。该归一化假设完整本底的相对变化服从两流表示比，但不推导未映射瞬发正电子残余的独立目标率。$N_{\mathrm{eff},w}$ 只描述转移权重集中度，既不是正式总本底有效样本量，也不是置信区间。}
\label{tab:environment_screening}
\scriptsize
\resizebox{\linewidth}{!}{%
\begin{tabular}{llrrrr}
\toprule
情景 & 源场代理 &
$ {\cal R}^{(2)}_e$ &
$F_{3\sigma,e}^{\mathrm{screen}}/F_{3\sigma,\mathrm B}$ &
$F_{3\sigma,e}^{\mathrm{screen}}$ / $\phcms$ &
$N_{\mathrm{eff},w}$ \\
\midrule
38 km 气球 & 正式 81 节点轨迹 &
1.0000 & 1.0000 & $2.4357\times10^{-5}$ & -- \\
530 km LEO & 近赤道、非 SAA 静态源场 &
0.5627 & 0.4830 & $1.1764\times10^{-5}$ & 44.27 \\
暴露月面 & REDMoon Apollo-17 月壤代理 &
23.0228 & 3.0894 & $7.5247\times10^{-5}$ & 5.23 \\
日--地 L2（2014） & 太阳极大、GCR 低通量 &
5.9460 & 1.5700 & $3.8241\times10^{-5}$ & 5.73 \\
日--地 L2（2009） & 太阳极小、GCR 高通量 &
11.9177 & 2.2227 & $5.4139\times10^{-5}$ & 5.05 \\
\bottomrule
\end{tabular}}
\end{table}

日--地 L2 两个端点的排序与质子谱一致：2014 年较强的太阳调制降低 GCR 驱动的延迟响应，2009 年较弱的调制则使其升高。筛查范围为 $(3.824$--$5.414)\times10^{-5}\phcms$，即正式气球阈值的 1.570--2.223 倍，但绝对幅度的响应支撑较弱。4 条选后 $^{62}$Cu 事例共用同一活化键，而该键只由 1 条 1.366\,GeV 质子产生记录表示；相应 $N_{\mathrm{eff},w}$ 仅为 5.05--5.73。因此，这一响应揭示了物理上可解释的质子活化风险，而不是已经收敛的 L2 本底。

各环境缺失的物理项并不相同。LEO 代理未包含轨道倾角和地磁截止的时间变化、俘获粒子与 SAA、离开 SAA 后的衰变、地球边缘历史、航天器质量、姿态和占空比。月面代理未包含当地月壤与地形、着陆器或基地、能源系统、太阳高能粒子和重离子。日--地 L2 代理未包含太阳高能粒子、重离子、方向性软质子、航天器结构与姿态以及轨道相关中断。三类目标环境均未完成匹配的“瞬发--活化--延迟”输运及自身公共时间事例折叠。因此，这些指数用于识别后续本底问题，不能用于排列完整任务构型的性能。

\FloatBarrier

\section{讨论}
\label{sec:discussion}

选后事例来源把质量模型 A 的本底诊断与质量模型 B 的性能比较连接起来。质量模型 A 的 TES 位于多级冷盘下方，该区域贡献带权选后延迟率的 32.18\%，而铜在独立材料投影中占 70.87\%。质量模型 B 将六层焦面横向置于 BGO 包围的烟囱内；选后冷端率由 $(1.21398\pm0.26869)\times10^{-2}$ 变为 $(1.40604\pm1.09289)\times10^{-4}\cps$，但后者仅由 7 条来源记录支撑。选后有效面积变化小于 0.25\%，而第 15 天直接本底、20 d 累计本底和选后延迟率分别降低 4.94、4.72 和 10.53 倍，$3\sigma$ 流量阈值改善 2.227 倍。

这一分析顺序与近期仪器本底研究一致。COSI 在详细气球载荷中分别输运瞬发与活化分量，并与飞行数据比较其能谱和时间行为 \cite{Gallego2025Balloon}；DIXE 同样从详细 TES 载荷出发，追问哪些入射粒子、产生区域和物理过程构成探测谱 \cite{Tian2026DIXE}。本工作中总清单与选后来源的区别尤其关键：主导选后组集中在局部 Cu 与近场结构，因此只依据总活度会遗漏真正的几何杠杆。

来源图也说明了为什么在冷端铜结构下方继续堆叠被动材料并不是高效解法。被动层可通过康普顿散射和光电吸收衰减 511 keV 光子，但要覆盖向下投影需要足够厚度和质量，还会与光学、低温服务开口竞争，并增加可被活化的材料 \cite{NISTXCOM}。这里需要区分两个过程：511 keV 光子低于核场单光子对产生阈值 $2m_ec^2=1.022$ MeV，本身不能发生电子--正电子对产生；但更高能的宽带光子和粒子级联可在新增材料中产生正电子，$\beta^+$ 活化也会产生正电子，随后湮没形成次级 511 keV 光子。质量模型 B 将 TES 置于活化铜直接投影之外，目的在于先削弱这种耦合，再由主动屏蔽处理剩余穿透事例。其紧凑 BGO 布局也提示可进一步节省主动屏蔽质量预算，但当前完整设计比较尚未单独量化这一节省。

选择谱同时分离了仪器各层的作用。主动 BGO 提供质量模型 B 的主要抑制，使科学窗直接率由 $2.38372$ 降至 $1.11379\times10^{-2}\cps$；最终拓扑/孔径检验留下 $1.07472\times10^{-2}\cps$，同时保留测量窗内 $98.249\%$ 的聚焦信号。最终率由宽带大气 $\gamma$ 本底、延迟活化和\revblue{七类非 $\gamma$ 瞬发粒子}分别贡献 $49.830\%$、$33.316\%$ 和 $16.854\%$。瞬发分量的中心值仅由 3 条选后记录支撑，因此更稳妥的解释是质量模型 A 的活化主导变为质量模型 B 的混合残余，而不是精确排列三类分量的大小。

该望远镜的观测角色与测量核球、银盘及更广银河发射的 INTEGRAL/SPI 和 COSI 互补 \cite{Siegert2016AA,Kierans2020,Siegert2020COSI,Yoneda2025,Karwin2023}。Laue 透镜--TES 概念以天空覆盖换取集中束流和亚 keV 科学窗，适合紧致 511 keV 源、未分辨中心分量及定点后随观测。表~\ref{tab:primary_sensitivity_zh} 量化的是持续在源、窄线的 20 d 参考情形。

\subsection{跨环境解释}

第~\ref{sec:environment_screening}~节的筛查不改变正式气球结果。其较稳健的信息是区分不同环境驱动：安静非 SAA LEO 代理降低两个可表示本底流的本底，暴露月面代理与瞬发 $\gamma$ 响应强耦合，而安静日--地 L2 与稀疏的质子活化响应耦合。这些趋势能否在任务构型中保留，取决于当前省略的平台材料、方向性辐射场、活化时间史和观测计划。特别是，LEO 仪器离开 SAA 并恢复科学观测后仍可能保留延迟本底；日--地 L2 即使没有外部行星中子源，载荷内部仍可由带电粒子产生次级中子。

\subsection{适用范围与下一步}
\label{subsec:limitations}

任务数值是探测器--低温系统几何及固定 $45^\circ$ 仰角、合成连续在源轨迹的统计投影。计算对各几何的主动 BGO 通道采用 50 keV 离线阈值，并采用探测器级拓扑检验、全部八族瞬发和全部八族活化。Laue 光学与气球平台本底不在当前探测器--低温系统预算内；辐射场、截面、标定、重建和指向系统学也不在传播统计不确定度内。

质量模型 B 的结果以完整构型为条件。相对质量模型 A，其入口几何、局部 W 结构、BGO 布局、主动通道和部分近场材料清单同时变化。近似相等的聚焦有效面积建立了信号可比性，却不能证明两模型对弥散光子的接受度相同。因此，现有数据不能把灵敏度变化单独归因于烟囱位移、BGO 布局或其他任一构件。

当前统计精度主要受选后输运事例数有限且权重不等限制。质量模型 B 最终瞬发分量和宽带大气 $\gamma$ 分量分别只有 3 条和 12 条记录，冷端延迟估计也只有 7 条来源记录。后续验证应优先增加这些分量的支撑，重复产生位置抽样以检验空间收敛，以实际飞行计划和可见性历史替换合成轨迹，并标定能量响应与 BGO 触发效率。随后可使用聚焦斑和控制区构建空间--能谱似然，同时剖面化本底扰动参数 \cite{Cowan2011}。这些工作将扩展当前“来源--几何”逻辑，并量化超出本次条件性设计比较的系统学。

\section{结论}
\label{sec:conclusions}

质量模型 A 的分析指出，多级冷端及其附近结构中的活化 Cu 是选后延迟本底的重要来源。这一来源图促使本文检验质量模型 B 的完整横向烟囱与主动 BGO 构型；后者保留质量模型 A 的 $99.75\%$ 选后有效面积。第 15 天直接选后本底由 $5.30571\times10^{-2}$ 降至 $1.07472\times10^{-2}\cps$；在 $\fref=2.4\times10^{-4}\phcms$ 下，质量模型 B 的信号率为 $2.35000\times10^{-3}\cps$。沿固定 $45^\circ$ 的 20 d、81 节点轨迹，质量模型 B 累积 4028.15 个信号和 18568.89 个本底计数，计数显著度为 29.561，对应的大气层顶 $3\sigma$ 阈值为 $(2.436\pm0.223)\times10^{-5}\phcms$；该阈值相对质量模型 A 改善 2.227 倍。这一改进是两套完整探测器--低温系统设计的条件比较，而不是单一构件的消融结果。

响应加权的两流扩展给出 530\,km 非 SAA LEO 代理 0.483、暴露月面代理 3.089，以及两个安静日--地 L2 调制端点 1.570--2.223 的筛查指数。转移未映射由 3 条记录支撑的瞬发正电子残余，也未包含目标环境结构、活化时间史和公共时间闭合。因此，这些数值用于确定环境匹配模拟的优先问题，而不是性能预报。

\section*{数据与软件可用性}
正式发表前，应通过带版本号的公共代码仓库或机构存档提供复现本文图表所需的几何和源定义、随机性约定、分析与制图脚本、LaTeX 源文件及派生验证产物。

\section*{声明}
\textbf{经费来源} 投稿前补充。

\textbf{利益冲突} 投稿前补充。

\textbf{作者贡献} 投稿前补充。

\textbf{伦理审批与参与同意} 不适用。

\begin{thebibliography}{99}

\bibitem{Prantzos2011}
N. Prantzos, C. Boehm, A.M. Bykov, R. Diehl, K. Ferriere, N. Guessoum, P. Jean, J. Knoedlseder, A. Marcowith, I.V. Moskalenko, A. Strong, G. Weidenspointner, The 511 keV emission from positron annihilation in the Galaxy, Rev. Mod. Phys. 83 (2011) 1001--1056, doi:10.1103/RevModPhys.83.1001.

\bibitem{Jean2006}
P. Jean, J. Knoedlseder, V. Lonjou, M. Allain, J.-P. Roques, G.K. Skinner, B.J. Teegarden, G. Vedrenne, P. von Ballmoos, B. Cordier, P. Caraveo, R. Diehl, P. Durouchoux, P. Leleux, R. Schanne, G. Weidenspointner, Spectral analysis of the Galactic $e^+e^-$ annihilation emission, Astron. Astrophys. 445 (2006) 579--589, doi:10.1051/0004-6361:20053765.

\bibitem{Kierans2020}
C.A. Kierans, S.E. Boggs, A. Zoglauer, A.W. Lowell, C. Sleator, J. Beechert, T.J. Brandt, P. Jean, H. Lazar, J. Roberts, T. Siegert, J.A. Tomsick, P. von Ballmoos, Detection of the 511 keV Galactic positron annihilation line with COSI, Astrophys. J. 895 (2020) 44, doi:10.3847/1538-4357/ab89a9.

\bibitem{Siegert2016AA}
T. Siegert, R. Diehl, G. Khachatryan, M.G.H. Krause, F. Guglielmetti, J. Greiner, A.W. Strong, X. Zhang, Gamma-ray spectroscopy of positron annihilation in the Milky Way, Astron. Astrophys. 586 (2016) A84, doi:10.1051/0004-6361/201527510.

\bibitem{Siegert2016V404}
T. Siegert, R. Diehl, J. Greiner, M.G.H. Krause, A.M. Beloborodov, M. Cadolle Bel, F. Guglielmetti, J. Rodriguez, A.W. Strong, X. Zhang, Positron annihilation signatures associated with the outburst of the microquasar V404 Cygni, Nature 531 (2016) 341--343, doi:10.1038/nature16978.

\bibitem{Knodlseder2005AllSky}
J. Knoedlseder, P. Jean, V. Lonjou, G. Weidenspointner, N. Guessoum, R. Diehl, W. Gillard, M.J. Harris, G.K. Skinner, P. von Ballmoos, G. Vedrenne, P. Caraveo, B. Cordier, V. Schoenfelder, B.J. Teegarden, The all-sky distribution of 511 keV electron-positron annihilation emission, Astron. Astrophys. 441 (2005) 513--532, doi:10.1051/0004-6361:20042063.

\bibitem{Yoneda2025}
H. Yoneda, T. Siegert, S. Mittal, Imaging the positron annihilation line with 20 years of INTEGRAL/SPI observations, arXiv:2509.01066, 2025.

\bibitem{Siegert2020COSI}
T. Siegert, S.E. Boggs, J.A. Tomsick, A.C. Zoglauer, C.A. Kierans, C.C. Sleator, J. Beechert, T.J. Brandt, P. Jean, H. Lazar, A.W. Lowell, J.M. Roberts, P. von Ballmoos, Imaging the 511 keV positron annihilation sky with COSI, Astrophys. J. 897 (2020) 45, doi:10.3847/1538-4357/ab9607.

\bibitem{Siegert2019Kinematics}
T. Siegert, R. Diehl, G. Khachatryan, M.G.H. Krause, F. Guglielmetti, J. Greiner, A.W. Strong, X. Zhang, Constraints on positron annihilation kinematics in the inner Galaxy, Astron. Astrophys. 627 (2019) A126, doi:10.1051/0004-6361/201833856.

\bibitem{DeCesare2011}
G. De Cesare, Searching for the 511 keV annihilation line from galactic compact objects with the IBIS gamma ray telescope, Astron. Astrophys. 531 (2011) A56, doi:10.1051/0004-6361/201116516.

\bibitem{Barriere2009Laue}
N.M. Barriere, L. Natalucci, N. Abrosimov, P. von Ballmoos, P. Bastie, P. Courtois, M. Jentschel, J. Knodlseder, J. Rousselle, P. Ubertini, Soft gamma-ray optics: new Laue lens design and performance estimates, arXiv:0910.0488, 2009.

\bibitem{Barriere2011LauePrototype}
N.M. Barriere, J.A. Tomsick, S.E. Boggs, A. Lowell, P. von Ballmoos, Developing a second generation Laue lens prototype: high reflectivity crystals and accurate assembly, arXiv:1111.6700, 2011.

\bibitem{Weidenspointner2006MAX}
G. Weidenspointner, C.B. Wunderer, N. Barriere, A. Zoglauer, P. von Ballmoos, Monte Carlo study of detector concepts for the MAX Laue Lens gamma-ray telescope, arXiv:astro-ph/0603152, 2006.

\bibitem{Shirazi2023}
F. Shirazi, E. Gau, M.A. Hossen, D. Becker, D. Schmidt, D. Swetz, D. Bennett, D. Braun, F. Kislat, J. Gard, J. Mates, J. Weber, N.R. Cavero, S. Chun, L. Lisalda, A. West, B. Dev, F. Ferrer, R. Bose, J. Ullom, H. Krawczynski, The 511-CAM Mission: A pointed 511 keV gamma-ray telescope with a focal plane detector made of stacked transition edge sensor microcalorimeter arrays, J. Astron. Telesc. Instrum. Syst. 9 (2023) 024006, doi:10.1117/1.JATIS.9.2.024006.

\bibitem{Zhang2022TESGamma}
S. Zhang, J.-K. Xia, T. Sun, et al., Transition edge sensor based detector: from X-ray to gamma-ray, arXiv:2204.00010, 2022.

\bibitem{Hu2026}
K. Hu, M. Fritts, D. Becker, D. Schmidt, F. Zhou, A. Dasgupta, F. Kislat, M. Keller, S. Chun, A.G. Detoito, E. Gau, X. Jin, D. Bennett, D. Braun, J. Gard, J.A.B. Mates, J. Nobles, D. Radomski, N. Rodriguez Cavero, R. Snodgrass, D. Swetz, J. Ullom, S. Vengalil Menon, J. Weber, K. Wimalasena, H. Krawczynski, Design of a mission to measure the shape and substructure of the 511 keV gamma-ray line from the center of the Milky Way, J. Astron. Telesc. Instrum. Syst. 12 (2026) 025002, doi:10.1117/1.JATIS.12.2.025002.

\bibitem{Gallego2025Balloon}
S. Gallego, U. Oberlack, J. Lommler, C.M. Karwin, A. Zoglauer, P. Jean, P. von Ballmoos, C. Kierans, C.C. Sleator, J.A. Tomsick, S.E. Boggs, Bottom-up background simulations of the 2016 COSI balloon flight, Astrophys. J. 986 (2025) 116, doi:10.3847/1538-4357/add6a0.

\bibitem{Tian2026DIXE}
R. Tian, J. Mao, J. Liu, H. Jin, W. Cui, Simulation of non-X-ray background for the DIffuse X-ray Explorer (DIXE) mission, arXiv:2604.13569, 2026.

\bibitem{Caputo2017}
R. Caputo, F. Kislat, J. Racusin, on behalf of the AMEGO Team, AMEGO: Simulations of the instrument performance, Proc. 35th International Cosmic Ray Conference, PoS(ICRC2017) 783 (2018), doi:10.22323/1.301.0783.

\bibitem{Gottardi2021TESReview}
L. Gottardi, K. Nagayoshi, A review of X-ray microcalorimeters based on superconducting transition edge sensors for astrophysics and particle physics, Applied Sciences 11 (2021) 3793, doi:10.3390/app11093793.

\bibitem{Xie2026AlMn}
L. Xie, Y. Zhang, Z. Li, Z. Liu, S. Shu, J. Zhou, X. Li, H. Li, H. Gao, Y. Gu, X. Lu, Y. Zhao, C. Liu, Beyond one-thousandth energy resolution with an AlMn TES detector, arXiv:2602.11728, 2026.

\bibitem{Vavagiakis2017AlMnMagnetic}
E.M. Vavagiakis, S.W. Henderson, K. Zheng, H.-M. Cho, N.F. Cothard, B. Dober, S.M. Duff, P.A. Gallardo, G. Hilton, J. Hubmayr, K.D. Irwin, B.J. Koopman, D. Li, F. Nati, M.D. Niemack, C.D. Reintsema, S. Simon, J.R. Stevens, A. Suzuki, B. Westbrook, Magnetic sensitivity of AlMn TESes and shielding considerations for next-generation CMB surveys, arXiv:1710.08456, 2017.

\bibitem{NISTXCOM}
J.H. Hubbell, S.M. Seltzer, XCOM: Photon Cross Section Database, NIST Standard Reference Database 8 (XGAM), National Institute of Standards and Technology, doi:10.18434/T48G6X.

\bibitem{Sturner2003}
S.J. Sturner, C.R. Shrader, G. Weidenspointner, B.J. Teegarden, D. Attie, B. Cordier, R. Diehl, C. Ferguson, P. Jean, A. von Kienlin, P. Paul, F. Sanchez, S. Schanne, P. Sizun, G. Skinner, C.B. Wunderer, Monte Carlo simulations and generation of the SPI response, Astron. Astrophys. 411 (2003) L81--L84, doi:10.1051/0004-6361:20031171.

\bibitem{Agostinelli2003}
S. Agostinelli, J. Allison, K. Amako, et al., Geant4---a simulation toolkit, Nucl. Instrum. Methods Phys. Res. A 506 (2003) 250--303, doi:10.1016/S0168-9002(03)01368-8.

\bibitem{Allison2016}
J. Allison, K. Amako, J. Apostolakis, et al., Recent developments in Geant4, Nucl. Instrum. Methods Phys. Res. A 835 (2016) 186--225, doi:10.1016/j.nima.2016.06.125.

\bibitem{SanchezDelRio2011XOP}
M. Sanchez del Rio, R.J. Dejus, XOP v2.4: recent developments of the x-ray optics software toolkit, Proc. SPIE 8141 (2011) 814115, doi:10.1117/12.893911.

\bibitem{Guan2023BraggReflection}
F. Guan, M. Asai, D.A. Bartkoski, M. Kleckner, Z. Harel, M. Salehpour, Adding the X-ray Bragg reflection physical process in crystal to the Geant4 Monte Carlo simulation toolkit, part I: reflection from a crystal slab, Precision Radiation Oncology 7(1) (2023) 59--66.

\bibitem{Reiazi2025G4BraggReflection}
R. Reiazi, D. Lee, D. Bartkoski, M. Kleckner, M. Salehpour, G4BraggReflection for accurate modeling of Bragg reflection in perfect and mosaic crystals, J. Phys. D: Appl. Phys. 58 (2025) 295102, doi:10.1088/1361-6463/aded1d.

\bibitem{Zoglauer2006}
A. Zoglauer, R. Andritschke, F. Schopper, MEGAlib: The Medium Energy Gamma-ray Astronomy Library, New Astron. Rev. 50 (2006) 629--632, doi:10.1016/j.newar.2006.06.049.

\bibitem{Sato2015EXPACS}
T. Sato, Analytical model for estimating terrestrial cosmic ray fluxes nearly anytime and anywhere in the world: extension of PARMA/EXPACS, PLoS ONE 10 (2015) e0144679, doi:10.1371/journal.pone.0144679.

\bibitem{Sato2016PARMA4}
T. Sato, Analytical model for estimating the zenith angle dependence of terrestrial cosmic ray fluxes, PLoS ONE 11 (2016) e0160390, doi:10.1371/journal.pone.0160390.

\bibitem{Kondev2021NUBASE}
F.G. Kondev, M. Wang, W.J. Huang, S. Naimi, G. Audi, The NUBASE2020 evaluation of nuclear physics properties, Chinese Phys. C 45 (2021) 030001, doi:10.1088/1674-1137/abddae.

\bibitem{Kierans2017}
C.A. Kierans, S.E. Boggs, J.-L. Chiu, A. Lowell, C. Sleator, J.A. Tomsick, A. Zoglauer, M. Amman, H.-K. Chang, C.-H. Tseng, C.-Y. Yang, C.-H. Lin, P. Jean, P. von Ballmoos, The 2016 Super Pressure Balloon flight of the Compton Spectrometer and Imager, arXiv:1701.05558, 2017.

\bibitem{Karwin2023}
C.M. Karwin, T. Siegert, J. Beechert, J.A. Tomsick, T.A. Porter, M. Negro, C. Kierans, M. Ajello, I. Martinez Castellanos, A. Shih, A. Zoglauer, S.E. Boggs, Probing the Galactic diffuse continuum emission with COSI, arXiv:2310.12206, 2023.

\bibitem{Cumani2019LEO}
P. Cumani, M. Hernanz, J. Kiener, V. Tatischeff, A. Zoglauer,
Background for a gamma-ray satellite on a low-Earth orbit,
Exp. Astron. 47 (2019) 273--302,
doi:10.1007/s10686-019-09624-0.

\bibitem{Dobynde2021REDMoon}
M.I. Dobynde, J. Guo,
Radiation environment at the surface and subsurface of the Moon:
model development and validation,
J. Geophys. Res. Planets 126 (2021) e2021JE006930,
doi:10.1029/2021JE006930.

\bibitem{Lotti2021Athena}
S. Lotti et al.,
Review of the particle background of the Athena X-IFU instrument,
Astrophys. J. 909 (2021) 111,
doi:10.3847/1538-4357/abd94c.

\bibitem{Usoskin2005}
I.G. Usoskin, K. Alanko-Huotari, G.A. Kovaltsov, K. Mursula,
Heliospheric modulation of cosmic rays: monthly reconstruction for 1951--2004,
J. Geophys. Res. Space Phys. 110 (2005) A12108,
doi:10.1029/2005JA011250.

\bibitem{Kuznetsov2017}
N.V. Kuznetsov, H. Popova, M.I. Panasyuk,
Empirical model of long-time variations of galactic cosmic ray particle fluxes,
J. Geophys. Res. Space Phys. 122 (2017) 1463--1472,
doi:10.1002/2016JA022920.

\bibitem{Cowan2011}
G. Cowan, K. Cranmer, E. Gross, O. Vitells, Asymptotic formulae for likelihood-based tests of new physics, Eur. Phys. J. C 71 (2011) 1554, doi:10.1140/epjc/s10052-011-1554-0.

\end{thebibliography}
\end{document}
